\documentclass[11pt]{article}
\usepackage[margin=1in]{geometry}
\usepackage{amsmath,amssymb,amsthm}
\usepackage{booktabs}
\usepackage{hyperref}
\usepackage{tikz}
\usepackage{xcolor}
\usepackage{listings}
\usetikzlibrary{arrows.meta,positioning,fit,backgrounds,shapes.geometric,calc}

% Long \texttt{} file names are unbreakable; allow a little extra interword
% stretch so they do not overflow the text block.
\setlength{\emergencystretch}{3em}

\newtheorem{theorem}{Theorem}
\newtheorem{proposition}{Proposition}
\newtheorem{lemma}{Lemma}
\newtheorem{corollary}{Corollary}
\theoremstyle{definition}
\newtheorem{definition}{Definition}
\newtheorem{example}{Example}
\newtheorem{remark}{Remark}

\newcommand{\lo}{\underline}
\newcommand{\hi}{\overline}
\newcommand{\indep}{\perp\!\!\!\perp}
\newcommand{\DSE}{\mathcal{D}_{\mathrm{SE}}}
\newcommand{\DLCN}{\mathcal{D}_{\mathrm{LCN}}}
\newcommand{\ch}{\mathrm{ch}}
\newcommand{\pa}{\mathrm{pa}}
\newcommand{\atoms}{\mathrm{atoms}}
\newcommand{\restr}{\!\restriction\!}
\newcommand{\readoff}{\rho}
\newcommand{\ext}{\mathrm{ext}}

\title{Exactness of Marginal Inference in Logical Credal Networks\\
\large A consolidated reference: precise statements, complete proofs,
and a verified example for every result}
\author{IBM Research}
\date{}

\begin{document}
\maketitle

\begin{abstract}
This document is the single reference for \emph{when} the LCN library's marginal-inference engines
return exact bounds, and for \emph{what} they bound when they do not. It fixes one framework --- the
LCN distribution set $\DLCN$, the strong extension $\DSE\supseteq\DLCN$ of the compiled credal
network, and a bound-character taxonomy that always names its reference set --- and then places each
engine in it. Every proposition and theorem carries a \emph{complete} proof rather than a sketch, and
every technical result is accompanied by a concrete example whose numbers were produced by running
the shipped engines.

The central definition is \emph{family-realizability} (Definition~\ref{def:realizable}), stated here
for general chain-graph LCNs, in which a family is a whole undirected chain component conditioned on
its parents. Around it we prove: the inclusion $\DLCN\subseteq\DSE$ (Proposition~\ref{prop:incl});
a sufficient condition for equality (Proposition~\ref{prop:suff}) that, notably, does \emph{not}
require single-connectedness; the exactness of Credal~VE over $\DSE$ for unconditional marginals
(Theorem~\ref{thm:cve}); Interval~BP's outer guarantee and its exactness on trees
(Theorem~\ref{thm:ibp}); ApproxLP's inner guarantee over $\DSE$ (Theorem~\ref{thm:alp}); the
CredalJT exactness theorem at treewidth cost under hypotheses H1--H4 (Theorem~\ref{thm:d5}); and,
for reference, ARIEL's message-soundness lemma together with the precise form of its exactness
claim (Section~\ref{sec:ariel}).

Three corrections to the companion notes are folded in, each established by measurement rather than
argument. Credal~CTE is \emph{not} exact at $\epsilon{=}0$: on \texttt{d4\_biting} it returns
$P(B)=[0.05,0.65]$ strictly inside the exact $[0.04,0.72]$, so it is not even a valid outer bound.
The characterization ``$\DSE{=}\DLCN$ iff singly connected and fully realizable'' cannot be an
\emph{iff} over every atom --- \texttt{d4\_biting} has $\DSE\supsetneq\DLCN$ yet exact marginals.
And single-connectedness is not what governs Gap~B: on a \emph{loopy} realizable $k$-tree, Credal~VE
is exact on every atom. We also record where the machinery has genuine holes --- dominance pruning
interacts badly with ratio read-off, the conditional-query denominator floors are soundness
hypotheses, and \texttt{ExactInference}-global fails to certify several atoms of a five-atom
instance --- and we state these as caveats instead of concealing them.
\end{abstract}

\tableofcontents

\section{The framework}\label{sec:framework}

\subsection{LCNs and the LCN distribution set}
An LCN over binary atoms $\mathbf V=\{V_1,\dots,V_n\}$ is a finite set of probability
\emph{sentences}, each of Type-1, $\ell\le P(\varphi)\le u$, or Type-2,
$\ell\le P(\varphi\mid\psi)\le u$, over propositional formulas, together with the conditional
independencies of the \emph{Local Markov Condition} (LMC) of its primal graph. A candidate joint
$p\in\Delta(2^{\mathbf V})$ over the $2^n$ truth assignments is a member of the \emph{LCN
distribution set} $\DLCN$ when it satisfies every sentence and every LMC equality.

\paragraph{Type-2 sentences: the cleared form.} A Type-2 sentence is a ratio constraint, and a ratio
is undefined when its conditioning event has probability zero. Every engine in the library therefore
implements the algebraically \emph{cleared} form
\begin{equation}\label{eq:cleared}
  \ell\cdot P(\psi)\;\le\;P(\varphi\wedge\psi)\;\le\;u\cdot P(\psi),
\end{equation}
which is total (defined for all $p$) and \emph{bilinear} in $p$. We define $\DLCN$ with
\eqref{eq:cleared}, not with the ratio, because that is the set the engines actually target.

\begin{remark}[The two readings differ]\label{rem:cleared}
On $\{p:P(\psi)=0\}$ the cleared form \eqref{eq:cleared} is satisfied vacuously, whereas the ratio
$P(\varphi\mid\psi)$ is undefined. Consequently the cleared reading admits joints the ratio reading
excludes. Defining $\DLCN$ by the ratio --- as earlier versions of this note did --- would compare
every engine against a set that none of them optimizes over. All results below are stated for the
cleared reading; they transfer to the ratio reading on the subset where every conditioning event has
positive probability.
\end{remark}

\paragraph{LMC equalities.} Each LMC assertion $(X\indep Y\mid S)$ imposes, for \emph{every} joint
configuration $\mathbf x,\mathbf y,\mathbf s$ of the atom blocks $X,Y,S$, the bilinear equality
\begin{equation}\label{eq:lmc}
  P(\mathbf x,\mathbf y,\mathbf s)\,P(\mathbf s)=P(\mathbf x,\mathbf s)\,P(\mathbf y,\mathbf s),
\end{equation}
i.e.\ $X\indep Y\mid S$ under $p$. These are built by
\texttt{lmc\_constraint\_groups\_vec}
(\texttt{lcn/inference/utils/\allowbreak common.py}). Writing $g_s$ for the functional a sentence $s$ constrains
and $h_\iota(p)=0$ for the equalities of assertion $\iota$,
\begin{equation}\label{eq:dlcn}
  \DLCN=\bigl\{\,p\in\Delta(2^{\mathbf V})\;:\;
    \ell_s\le g_s(p)\le u_s\ \forall s,\qquad h_\iota(p)=0\ \forall \iota\,\bigr\},
\end{equation}
where a Type-1 $g_s(p)=P(\varphi)=\mathbf 1_\varphi^\top p$ is linear and a Type-2 sentence
contributes the two rows of \eqref{eq:cleared}. The marginal task for a query atom $a$ is
\[
  \lo P(a)=\min_{p\in\DLCN}P(a),\qquad \hi P(a)=\max_{p\in\DLCN}P(a),
\]
and analogously for a conditional $P(a\mid e)$ under evidence $e$.

\subsection{Locality}
Write $p\restr A$ for the marginal of $p$ on an atom set $A$. The following lemma is used by every
proof in this document; in the companion notes it appears as an unproved remark.

\begin{lemma}[Locality]\label{lem:locality}
Every constraint in \eqref{eq:dlcn} depends on $p$ only through its marginal on a fixed atom set,
its \emph{scope}: a sentence $s$ through $p\restr\atoms(s)$, and an LMC assertion $(X\indep Y\mid S)$
through $p\restr(X\cup Y\cup S)$. Consequently, if $p$ and $p'$ satisfy
$p\restr A=p'\restr A$ for the scope $A$ of a constraint $c$, then $c$ holds for $p$ iff it holds
for $p'$.
\end{lemma}

\begin{proof}
Let $c$ be a constraint with scope $A$. There are three cases.

\emph{Type-1.} $P(\varphi)=\sum_{\omega\models\varphi}p(\omega)$. Since
$\atoms(\varphi)\subseteq A$, the truth value of $\varphi$ at $\omega$ is determined by
$\omega\restr A$; grouping the sum by the value of $\omega\restr A$ gives
$P(\varphi)=\sum_{\alpha\models\varphi}(p\restr A)(\alpha)$, where $\alpha$ ranges over assignments
of $A$. This is a function of $p\restr A$ alone.

\emph{Type-2.} By the previous case both $P(\varphi\wedge\psi)$ and $P(\psi)$ are functions of
$p\restr A$ with $A\supseteq\atoms(\varphi)\cup\atoms(\psi)$; the rows \eqref{eq:cleared} are
therefore functions of $p\restr A$.

\emph{LMC.} Each factor in \eqref{eq:lmc} is a probability of a conjunction of literals over
$X\cup Y\cup S$, hence by the Type-1 case a function of $p\restr(X\cup Y\cup S)$; a product of such
functions is again one.

In each case the constraint's truth value is determined by $p\restr A$, which is precisely the
stated consequence.
\end{proof}

\subsection{Chain graphs, families, and the credal network}
The LCN's structure graph is simplified by collapsing each connected component of its undirected
part into a meta-node (\texttt{LCN.simplify\_structure\_graph}). When the result is a directed
acyclic graph the LCN is a \emph{chain graph}, and the joint factorizes over \emph{families}
\begin{equation}\label{eq:factor}
  P(\mathbf V)=\prod_v P(\ch_v\mid\pa_v),
\end{equation}
one per node $v$ of the simplified graph, where $\ch_v$ is the (possibly multi-atom) chain component
at $v$ and $\pa_v$ its parents. The family's \emph{scope} is $\ch_v\cup\pa_v$. Root families have
$\pa_v=\varnothing$, so their factor is the marginal $P(\ch_v)$.

\begin{example}[$\ch_v$ is genuinely a block]\label{ex:blocks}
Multi-atom chain components are the rule, not an edge case. Running
\texttt{LCN.process\_chain\_graph}:
\begin{itemize}\itemsep1pt
  \item \texttt{examples/alarm.lcn} yields families $A\mid B,E$; $B$; $E$; and
    $\boxed{C\text{-}D\mid A}$ --- a two-atom component $\ch_v=\{C,D\}$, because \texttt{s4}
    couples $C$ and $D$ undirectedly.
  \item \texttt{examples/smokers.lcn} yields $\boxed{F_1F_2F_3\mid\varnothing}$ and
    $\boxed{S_1S_2S_3\mid F_1F_2F_3}$ --- two \emph{three}-atom components --- besides
    $C_i\mid S_1S_2S_3$.
\end{itemize}
Any definition of realizability phrased for a single child atom cannot even be stated on these
instances; Definition~\ref{def:realizable} is therefore given for blocks.
\end{example}

\paragraph{Local credal sets and the strong extension.} For each family the \texttt{cn/} pipeline
(\texttt{local\_credal\_sets.py} with \texttt{method="linear"}) computes a \emph{local credal
set} $K_v$: the set of conditional tables $P(\ch_v\mid\pa_v)$ attainable by \emph{some} joint satisfying
the sentences whose atom scope lies inside the family scope. Operationally it is a linear-fractional
program $P(\ch_v,\pa_v)/P(\pa_v)$ solved by Charnes--Cooper, in which \emph{the parents'
distribution $P(\pa_v)$ is a free variable}. The credal-network engines bound the \emph{strong
extension}
\begin{equation}\label{eq:se}
  \DSE=\Bigl\{\textstyle\prod_v P_v\;:\;P_v\in K_{v,\pi}\text{ chosen independently for each
  parent configuration }\pi\Bigr\}.
\end{equation}

\begin{remark}[Row-wise versus table-wise, a third looseness source]\label{rem:rowwise}
Equation~\eqref{eq:se} is the \emph{row-independent} product $\prod_v\prod_\pi K_{v,\pi}$: the code
(\texttt{cve.py}, via \texttt{itertools.product} over parent configurations) chooses one conditional
row per parent configuration independently. This is in general \emph{strictly larger} than
$\{$tables lying in $K_v\}$, since $K_v$ need not be a Cartesian product of its rows. The companion
notes use both readings interchangeably; we fix the row-wise reading of \eqref{eq:se} throughout,
because it is the one implemented. Note that this is a source of looseness distinct from Gap~A and
Gap~B below, and unnamed by that taxonomy.
\end{remark}

\begin{remark}[The vertex machinery needs \texttt{method="linear"}]\label{rem:d1convex}
Under \texttt{method="linear"} each $K_{v,\pi}$ is a polytope (the Charnes--Cooper image of a
polyhedron), so it has finitely many extreme points and $\DSE$ is a product of polytopes --- the
hypothesis every vertex-enumeration argument below needs. Under scheme~D1
(\texttt{method="linear-tight"}) the local program is additionally cut by \emph{bilinear} LMC
equalities, so $K_{v,\pi}$ need not be convex; vertex enumeration then explores a subset of the
convex hull, and an engine that is exact over $\DSE$ for \texttt{"linear"} may become \emph{inner}
with respect to the D1 strong extension. All exactness results below are stated for
\texttt{method="linear"}.
\end{remark}

\subsection{Bound character, relative to a named reference set}
\begin{definition}[Bound character]\label{def:char}
Let $\mathcal R$ be a reference set (here $\DLCN$ or $\DSE$) with true bounds
$[\lo q^\star,\hi q^\star]$ for the query at hand. An engine's interval $[\lo q,\hi q]$ is
\emph{exact w.r.t.\ $\mathcal R$} if $[\lo q,\hi q]=[\lo q^\star,\hi q^\star]$; a valid \emph{outer}
bound if $[\lo q,\hi q]\supseteq[\lo q^\star,\hi q^\star]$; an \emph{inner} bound if
$[\lo q,\hi q]\subseteq[\lo q^\star,\hi q^\star]$; and \emph{unreliable} if neither containment is
guaranteed.
\end{definition}

Naming $\mathcal R$ is not pedantry: an engine can be inner w.r.t.\ $\DSE$ and simultaneously outer
w.r.t.\ $\DLCN$, since $\DLCN\subseteq\DSE$. ApploxLP is exactly such a case
(Example~\ref{ex:alpouter}), and the unqualified claim ``ApproxLP is not a valid outer bound'' found
in earlier versions of this note is false as stated.

\subsection{The two looseness gaps}\label{sec:gaps}
\begin{itemize}
  \item \textbf{Gap A (local-set widening).} Each $K_v$ is computed per family in isolation, so it
    may admit conditionals that no global $p\in\DLCN$ realizes: it ignores sentences on \emph{other}
    families and the cross-family LMC equalities. Formally
    $K_v^{\mathrm{true}}\subseteq K_v^{\mathrm{local}}$, sometimes strictly.
  \item \textbf{Gap B (strong-extension widening).} Even with $K_v=K_v^{\mathrm{true}}$, the free
    product of \eqref{eq:se} lets each family pick its row independently, which can violate a
    sentence whose scope spans several families.
\end{itemize}
Schemes D1 (\texttt{method="linear-tight"}), D2 (\texttt{merge\_budget}) and D3
(\texttt{solver="scip"}) attack Gap~A at build time and are free for every credal-network engine;
D4 (\texttt{coupling="cross-family"}) and D5 (\texttt{CredalJT}) attack Gap~B at propagation time.
See \texttt{tighter\_approximation.tex} for the taxonomy.

\section{Family-realizable sentences}\label{sec:realizable}

Which sentences can a \emph{single} family's local set enforce? A family's degrees of freedom are the
conditional probabilities $P(\ch_v{=}c\mid\pa_v{=}\pi)$; the parents' marginal $P(\pa_v)$ is not one
of them --- it is determined elsewhere in the network, and is a free variable during the local solve.
A sentence that constrains a quantity depending on $P(\pa_v)$ therefore cannot be enforced locally,
and the free product \eqref{eq:se} may violate it. This is the mechanism behind Gap~B, and the
following definition isolates it.

\begin{definition}[Family-realizable sentence]\label{def:realizable}
Let $\mathcal L$ be a chain-graph LCN with families $\{(\ch_v,\pa_v)\}$, and let $s$ be a sentence
with functional $g_s$. Say $s$ is \emph{family-realizable in family $v$} if $\atoms(s)\subseteq
\ch_v\cup\pa_v$ and the value $g_s(p)$ is determined by the conditional table
$P(\ch_v\mid\pa_v)$ \emph{alone} --- i.e.\ for any two joints $p,p'$ that induce the same conditional
table at $v$, $g_s(p)=g_s(p')$, whatever the parents' marginals $P(\pa_v)$ and $P'(\pa_v)$ may be.
Concretely, $s$ is family-realizable in $v$ if and only if it has one of the two forms:
\begin{enumerate}
  \item[\textup{(R1)}] \emph{Type-1 on a root block:} $s$ is $\ell\le P(\varphi)\le u$ with
    $\atoms(\varphi)\subseteq\ch_v$ and $\pa_v=\varnothing$. Then $P(\varphi)$ is a function of
    $P(\ch_v)$, which \emph{is} the family's table.
  \item[\textup{(R2)}] \emph{Type-2 pinning one full parent configuration:} $s$ is
    $\ell\le P(\varphi\mid\psi)\le u$ with $\atoms(\varphi)\subseteq\ch_v$,
    $\atoms(\psi)\subseteq\pa_v$, and $\psi$ satisfied by \emph{exactly one} assignment $\pi^\star$
    of $\pa_v$. Then $P(\varphi\mid\psi)=P(\varphi\mid\pa_v{=}\pi^\star)$ is a single row of the
    table.
\end{enumerate}
A sentence is \emph{non-realizable} if it is family-realizable in no family. The LCN is \emph{fully
realizable} if every sentence is family-realizable in some family.
\end{definition}

Note the quantifier: ``in \emph{some} family''. The library attaches a sentence to every family whose
scope contains its atoms (\texttt{process\_chain\_graph}), so one sentence is typically attached to
several families while being realizable in only one of them --- see Example~\ref{ex:fr5}.

The next lemma is what makes Definition~\ref{def:realizable} a \emph{criterion} rather than a
description: it shows the two forms (R1)--(R2) exhaust the $P(\pa_v)$-invariant sentences, so
realizability can be decided syntactically.

\begin{lemma}[Invariance criterion]\label{lem:invariance}
Let $v$ be a family with $\pa_v\ne\varnothing$ and let $s$ be a sentence with
$\atoms(s)\subseteq\ch_v\cup\pa_v$. Write $\pi$ for assignments of $\pa_v$, $w_\pi=P(\pa_v{=}\pi)$
for the parents' marginal, and $T(\cdot\mid\pi)=P(\ch_v{=}\cdot\mid\pa_v{=}\pi)$ for the table.
\begin{enumerate}
  \item[\textup{(i)}] If $s$ is Type-1, $\ell\le P(\varphi)\le u$, then
    $P(\varphi)=\sum_\pi w_\pi\,T(\varphi_\pi\mid\pi)$ where
    $T(\varphi_\pi\mid\pi):=\sum_{c\,:\,(c,\pi)\models\varphi}T(c\mid\pi)$. This is
    $P(\pa_v)$-invariant for all tables iff the map $\pi\mapsto T(\varphi_\pi\mid\pi)$ is constant
    for every table --- which fails whenever $|\{\pi\}|\ge2$ and $\varphi$ is non-trivial. Hence no
    Type-1 sentence is realizable in a non-root family.
  \item[\textup{(ii)}] If $s$ is Type-2 with $\atoms(\psi)\subseteq\pa_v$ and
    $\Pi_\psi:=\{\pi:\pi\models\psi\}$, then on $\{P(\psi)>0\}$
    \begin{equation}\label{eq:mixture}
      P(\varphi\mid\psi)=\sum_{\pi\in\Pi_\psi}\lambda_\pi\,T(\varphi_\pi\mid\pi),
      \qquad
      \lambda_\pi=\frac{w_\pi}{\sum_{\pi'\in\Pi_\psi}w_{\pi'}},
    \end{equation}
    a convex combination of table rows with weights $\lambda$ determined by $P(\pa_v)$. If
    $|\Pi_\psi|=1$ this is one row, hence invariant; if $|\Pi_\psi|\ge2$ it is invariant for all
    tables iff the rows agree, so $s$ is \emph{not} realizable.
\end{enumerate}
\end{lemma}

\begin{proof}
(i) Condition the Type-1 probability on the parent block. Since
$\atoms(\varphi)\subseteq\ch_v\cup\pa_v$, the truth value of $\varphi$ at a joint assignment is
determined by $(\ch_v,\pa_v)$, so by Lemma~\ref{lem:locality}
$P(\varphi)=\sum_\pi P(\pa_v{=}\pi)\,P(\varphi\mid\pa_v{=}\pi)=\sum_\pi w_\pi T(\varphi_\pi\mid\pi)$,
which is the stated formula. This is an affine function of the weight vector $w$ on the simplex, and
an affine function on a simplex is constant iff all its coefficients coincide; hence invariance for
\emph{every} table forces $\pi\mapsto T(\varphi_\pi\mid\pi)$ to be constant. Choose a table with
$T(\varphi_{\pi_1}\mid\pi_1)=1$ and $T(\varphi_{\pi_2}\mid\pi_2)=0$ for two distinct parent
configurations $\pi_1\ne\pi_2$ (possible whenever $\varphi$ is satisfiable and refutable within
$\ch_v$, i.e.\ non-trivial); then the coefficients differ and invariance fails. So $s$ is not
realizable in $v$.

(ii) With $\atoms(\psi)\subseteq\pa_v$, the event $\psi$ is a union of parent configurations,
$\{\psi\}=\bigcup_{\pi\in\Pi_\psi}\{\pa_v{=}\pi\}$, a disjoint union. Hence
$P(\psi)=\sum_{\pi\in\Pi_\psi}w_\pi$ and, by the computation in (i) restricted to $\Pi_\psi$,
$P(\varphi\wedge\psi)=\sum_{\pi\in\Pi_\psi}w_\pi T(\varphi_\pi\mid\pi)$. Dividing gives
\eqref{eq:mixture}, and $\lambda$ is a probability vector supported on $\Pi_\psi$.

If $|\Pi_\psi|=1$, say $\Pi_\psi=\{\pi^\star\}$, then $\lambda_{\pi^\star}=1$ and
$P(\varphi\mid\psi)=T(\varphi_{\pi^\star}\mid\pi^\star)$, a single row --- invariant, and this is
(R2). If $|\Pi_\psi|\ge2$, then \eqref{eq:mixture} is a non-degenerate convex combination: pick
$\pi_1\ne\pi_2$ in $\Pi_\psi$ and a table with $T(\varphi_{\pi_1}\mid\pi_1)=1$,
$T(\varphi_{\pi_2}\mid\pi_2)=0$. Taking $w$ concentrated near $\pi_1$ gives
$P(\varphi\mid\psi)\to1$, and near $\pi_2$ gives $\to0$, so the value genuinely depends on
$P(\pa_v)$ and $s$ is not realizable. Equivalently, $\partial P(\varphi\mid\psi)/\partial\lambda_\pi
=T(\varphi_\pi\mid\pi)$ is not constant across $\Pi_\psi$.
\end{proof}

\begin{remark}[Atom scope is not the criterion]\label{rem:scope}
Lemma~\ref{lem:invariance} shows family-realizability is strictly stronger than the atom-scope test
``$\atoms(s)$ fits inside a family scope''. Three canonical offenders are all atom-scope
family-local:
\begin{itemize}\itemsep1pt
  \item a bare marginal $P(x)$ on a \emph{non-root} child (Lemma~\ref{lem:invariance}(i));
  \item a child--parent joint $P(x\wedge y)$ with $y=\pa_x$, which equals $P(x\mid y)P(y)$;
  \item a Type-2 sentence whose $\psi$ is satisfied by two or more parent configurations --- e.g.\
    $P(x_2\mid x_0\oplus x_1)$ at a collider (Lemma~\ref{lem:invariance}(ii)). Note this includes a
    conjunction over a \emph{strict subset} of the parents, which pins a set of configurations rather
    than one.
\end{itemize}
\end{remark}

\subsection{A five-variable chain-graph example}\label{sec:fr5}

\begin{example}[Family-realizability on a five-atom chain graph]\label{ex:fr5}
The instance \texttt{ktree\_fr\_k3\_n5\_1.lcn} from \texttt{benchmarks/ktree\_small\_fr/}
($n{=}5$, a family-realizable $k$-tree with $k{=}3$):
\begin{align*}
  s_0:&\ 0.01\le P(x_3)\le0.399975 &
  s_4:&\ 0.673756\le P(\lnot x_0\mid x_3\wedge x_4\wedge\lnot x_1)\le1.0\\
  s_1:&\ 0.01\le P(x_4\mid x_3)\le0.442867 &
  s_5:&\ 0.748845\le P(\lnot x_3)\le0.954666\\
  s_2:&\ 0.638553\le P(x_2\mid\lnot x_3\wedge x_4)\le1.0 &
  s_6:&\ 0.141404\le P(x_3)\le0.247704\\
  s_3:&\ 0.01\le P(\lnot x_1\mid x_3\wedge x_4\wedge x_2)\le0.307066
\end{align*}
\texttt{LCN.process\_chain\_graph} reports five singleton-block families and a single LMC assertion:
\[
  x_3\mid\varnothing,\quad x_4\mid x_3,\quad x_2\mid x_3,x_4,\quad
  x_1\mid x_2,x_3,x_4,\quad x_0\mid x_1,x_3,x_4;
  \qquad (x_0\indep x_2\mid x_1,x_3,x_4).
\]
Table~\ref{tab:fr5} classifies every sentence. Two points are worth drawing out. First, the
\emph{attached} column (what \texttt{process\_chain\_graph} collects by the scope-subset test) is
much larger than the \emph{realizable-in} column: the root marginals $s_0,s_5,s_6$ are attached to
all five families but realizable only in the root family $x_3$, which is exactly why
Definition~\ref{def:realizable} quantifies existentially over families. Second, each of
$s_1,s_2,s_3,s_4$ conditions on a \emph{full} conjunction of signed literals over all of its
family's parents, so $|\Pi_\psi|=1$ and (R2) applies --- this is what the generator's
\texttt{full\_parents=True} flag guarantees for every \texttt{-fr} class.

\begin{table}[ht]\centering\small
\caption{Sentence classification for \texttt{ktree\_fr\_k3\_n5\_1.lcn}. Every sentence is
family-realizable in exactly one family, so the instance is fully realizable.}\label{tab:fr5}
\begin{tabular}{@{}llllc@{}}
\toprule
sentence & form & attached to families & realizable in & why \\
\midrule
$s_0,s_5,s_6$ & Type-1 $P(x_3)$, $P(\lnot x_3)$ & all five & $x_3\mid\varnothing$ & (R1) root block \\
$s_1$ & Type-2, $\psi=x_3$            & $x_4,x_2,x_1,x_0$ & $x_4\mid x_3$        & (R2) $|\Pi_\psi|{=}1$ \\
$s_2$ & Type-2, $\psi=\lnot x_3\wedge x_4$ & $x_2,x_1$    & $x_2\mid x_3,x_4$    & (R2) $|\Pi_\psi|{=}1$ \\
$s_3$ & Type-2, $\psi=x_3\wedge x_4\wedge x_2$ & $x_1$     & $x_1\mid x_2,x_3,x_4$& (R2) $|\Pi_\psi|{=}1$ \\
$s_4$ & Type-2, $\psi=x_3\wedge x_4\wedge\lnot x_1$ & $x_0$ & $x_0\mid x_1,x_3,x_4$& (R2) $|\Pi_\psi|{=}1$ \\
\bottomrule
\end{tabular}
\end{table}

\noindent\emph{Measured bounds.} \texttt{ExactInference} with \texttt{solver="global"} certifies all
ten solves (gap $0$, status \texttt{confirmed}), and \texttt{CredalJT} reproduces every endpoint:
\begin{center}\small
\begin{tabular}{@{}lccccc@{}}
\toprule
engine & $P(x_0)$ & $P(x_1)$ & $P(x_2)$ & $P(x_3)$ & $P(x_4)$ \\
\midrule
\texttt{ExactInference} (global) & $[0,1]$ & $[0,1]$ & $[0,1]$ & $[0.141404,0.247704]$ & $[0.001414,0.921219]$ \\
\texttt{CredalJT} (D5, SCIP)     & $[0,1]$ & $[0,1]$ & $[0,1]$ & $[0.141404,0.247704]$ & $[0.001414,0.921219]$ \\
\bottomrule
\end{tabular}
\end{center}
The vacuous $[0,1]$ on $x_0,x_1,x_2$ is genuine, not a solver artifact: \texttt{cjt\_exactness.tex}
exhibits explicit feasible witnesses attaining both endpoints (e.g.\ setting
$P(x_2\mid\lnot x_3\wedge x_4)=0$ makes $s_3$'s conditioning event unreachable, freeing $x_1$).
The junction tree has maximal cliques $\{x_0,x_1,x_3,x_4\}$ and $\{x_1,x_2,x_3,x_4\}$ --- treewidth
$3$ --- and the separator between them is exactly $\{x_1,x_3,x_4\}$, the conditioning set of the lone
LMC assertion; this is the certificate that hypotheses H1--H4 of Theorem~\ref{thm:d5} hold.
\end{example}

\begin{example}[The Type-2 half, isolated by a matched pair]\label{ex:collider}
To see that clause (R2)'s ``exactly one parent configuration'' is doing real work --- and not merely
excluding pathologies --- compare two LCNs that differ \emph{only} in the shape of $\psi$ at a
collider $x_0,x_1\to x_2$, with identical structure and identical numeric bounds:
\begin{center}\small
\texttt{P0}: $0.3\le P(x_0)\le0.5$;\quad
\texttt{P1}: $0.4\le P(x_1)\le0.6$;\quad
\texttt{C}: $0.2\le P(x_2\mid\psi)\le0.5$.
\end{center}
\begin{center}\small
\begin{tabular}{@{}llccc@{}}
\toprule
$\psi$ & $|\Pi_\psi|$ & realizable? & \texttt{ExactInference} (global) $P(x_2)$ & Credal~VE $P(x_2)$\\
\midrule
$x_0\oplus x_1$   & $2$ & \textbf{no}  & $[0.092,0.770]$ & $\mathbf{[0,1]}$ \emph{vacuous} \\
$x_0\wedge x_1$   & $1$ & yes          & $[0.024,0.940]$ & $\mathbf{[0.024,0.940]}$ \emph{exact} \\
\bottomrule
\end{tabular}
\end{center}
With $\psi=x_0\oplus x_1$ the sentence bounds a $P(\pa_{x_2})$-weighted mixture of two table rows
(Lemma~\ref{lem:invariance}(ii)); neither row is individually constrained, so every
$K_{x_2,\pi}$ is the full box $[0,1]$ and Credal~VE returns the vacuous interval, losing the bound
entirely. With $\psi=x_0\wedge x_1$ the same sentence pins the single row $\pi^\star=(1,1)$, clause
(R2) applies, and Credal~VE is exact. Both $P(x_0)$ and $P(x_1)$ are exact ($[0.3,0.5]$,
$[0.4,0.6]$) in both cases: the damage is local to the offending family. Note also that this is not
a solver or enumeration artifact --- \texttt{strong\_extension\_exactness.tex} verifies that neither
the global SCIP backend nor scheme~D1 recovers the bound, because the coupling is not representable
as a product of per-configuration credal sets.
\end{example}

\begin{example}[The Type-1 half: a marginal on a non-root child]\label{ex:2node}
The smallest witness. Take the two-atom chain $x_0\to x_1$ with
\begin{center}\small
\texttt{P0}: $0.3\le P(x_0)\le0.5$;\quad
\texttt{T0}: $0.2\le P(x_1\mid x_0)\le0.4$;\quad
\texttt{T1}: $0.6\le P(x_1\mid\lnot x_0)\le0.8$,
\end{center}
and optionally add \texttt{M1}: $0.45\le P(x_1)\le0.55$, a Type-1 marginal on the non-root child
$x_1$ --- non-realizable by Lemma~\ref{lem:invariance}(i). Measured:
\begin{center}\small
\begin{tabular}{@{}lcc@{}}
\toprule
engine & $P(x_1)$ without \texttt{M1} & $P(x_1)$ with \texttt{M1} \\
\midrule
\texttt{ExactInference} (global) & $[0.40,0.68]$ & $\mathbf{[0.45,0.55]}$ \\
\texttt{CredalJT} (D5)           & $[0.40,0.68]$ & $\mathbf{[0.45,0.55]}$ \\
Credal~VE                        & $[0.40,0.68]$ & $[0.40,0.68]$ \\
Interval~BP (\texttt{interval})  & $[0.40,0.68]$ & $[0.40,0.68]$ \\
ApproxLP                         & $[0.40,0.68]$ & $[0.40,0.68]$ \\
\bottomrule
\end{tabular}
\end{center}
Without \texttt{M1} the LCN is fully realizable and all five engines agree. With \texttt{M1} the
three strong-extension engines are unchanged --- the family $P(x_1\mid x_0)$ cannot enforce a
constraint on $P(x_1)$, which depends on $P(x_0)$ --- while the two $\DLCN$ engines tighten to the
sentence bound. $P(x_0)=[0.30,0.50]$ for every engine in both instances.
\end{example}

\section{The LCN set versus the strong extension}\label{sec:incl}

We now relate $\DLCN$ and $\DSE$. Three separate statements are needed, and it is important that they
are \emph{not} packaged as a single biconditional --- Remark~\ref{rem:nobiconditional} explains why.

\begin{proposition}[Inclusion]\label{prop:incl}
Let $\mathcal L$ be a chain-graph LCN. Then $\DLCN\subseteq\DSE$. Consequently, for every query,
$[\lo q_{\mathrm{LCN}},\hi q_{\mathrm{LCN}}]\subseteq[\lo q_{\mathrm{SE}},\hi q_{\mathrm{SE}}]$: any
engine that is exact or outer w.r.t.\ $\DSE$ is automatically a valid \emph{outer} bound w.r.t.\
$\DLCN$.
\end{proposition}

\begin{proof}
Let $p\in\DLCN$. Since $\mathcal L$ is a chain graph and $p$ satisfies every LMC equality, the
chain-graph factorization theorem (Lauritzen--Wermuth--Frydenberg; Lauritzen, \emph{Graphical
Models}, Thm.~3.34) gives the factorization \eqref{eq:factor}:
$p=\prod_v P_p(\ch_v\mid\pa_v)$, where $P_p(\cdot\mid\cdot)$ denotes the conditionals induced by $p$
(defined arbitrarily on parent configurations of zero probability, which do not affect the product).

Fix a family $v$ and a parent configuration $\pi$ with $P_p(\pa_v{=}\pi)>0$. The row
$P_p(\ch_v\mid\pa_v{=}\pi)$ is a conditional table attainable by the joint $p$, and $p$ satisfies
every sentence of $\mathcal L$ --- in particular every sentence whose atom scope lies inside
$\ch_v\cup\pa_v$. But $K_{v,\pi}$ was defined as the set of rows attainable by \emph{some} joint
satisfying exactly those in-scope sentences. Hence $P_p(\ch_v\mid\pa_v{=}\pi)\in K_{v,\pi}$: the
local set is a \emph{relaxation} of the projection, so it contains it. For a configuration with
$P_p(\pa_v{=}\pi)=0$ the row is unconstrained by $p$ and may be chosen to be any element of
$K_{v,\pi}$ (non-empty, since $p$ itself witnesses feasibility of the in-scope sentences) without
changing the product.

Thus $p$ is exhibited as a product $\prod_v\prod_\pi$ of rows drawn from the respective
$K_{v,\pi}$, which is precisely membership in $\DSE$ as defined in \eqref{eq:se}. The bound
inclusion follows because minimizing and maximizing a fixed functional over a larger set yields a
wider interval.
\end{proof}

\begin{proposition}[Sufficiency for equality]\label{prop:suff}
Let $\mathcal L$ be a chain-graph LCN with local sets built by \texttt{method="linear"}. Suppose
\begin{enumerate}
  \item[\textup{(A)}] \emph{no Gap~A:} for every family $v$ and parent configuration $\pi$,
    $K_{v,\pi}=K_{v,\pi}^{\mathrm{true}}$, the projection of $\DLCN$ onto that row; and
  \item[\textup{(B)}] \emph{full realizability:} every sentence of $\mathcal L$ is family-realizable
    in some family (Definition~\ref{def:realizable}).
\end{enumerate}
Then $\DSE=\DLCN$. \emph{Single-connectedness of the primal graph is not required.}
\end{proposition}

\begin{proof}
Proposition~\ref{prop:incl} gives $\DLCN\subseteq\DSE$, so only $\DSE\subseteq\DLCN$ needs proof.
Let $q=\prod_v\prod_\pi P_{v,\pi}$ with each $P_{v,\pi}\in K_{v,\pi}$ be an arbitrary element of
$\DSE$. We verify that $q$ satisfies both families of constraints in \eqref{eq:dlcn}.

\emph{LMC equalities.} By construction $q$ is a product of the form \eqref{eq:factor} over the
families of the chain graph. By the converse direction of the chain-graph factorization theorem
(Lauritzen, Thm.~3.34: factorization implies the global Markov property of the chain graph, and this
direction needs no positivity assumption), $q$ satisfies every conditional independence entailed by
the chain graph --- in particular every LMC assertion, since the LMC is by definition read off that
same graph. Hence $h_\iota(q)=0$ for all $\iota$. Note this step is where the ``loopiness contributes
nothing'' claim comes from: a product over the chain-graph families satisfies the LMC whether or not
the graph is singly connected.

\emph{Sentences.} Let $s$ be a sentence. By (B) it is family-realizable in some family $v$, so by
Definition~\ref{def:realizable} the value $g_s(q)$ is determined by the conditional table of $q$ at
$v$ alone. That table is the collection of rows $\{P_{v,\pi}\}_\pi$, and by (A) each
$P_{v,\pi}\in K_{v,\pi}^{\mathrm{true}}$, i.e.\ each row is realized by \emph{some} $p^{(\pi)}\in
\DLCN$. It remains to see that this forces $g_s(q)\in[\ell_s,u_s]$. Two cases, matching the two
clauses of Definition~\ref{def:realizable}.

If $s$ has form (R2), then $g_s(q)=P_{v,\pi^\star}(\varphi_{\pi^\star})$ is the \emph{single} row
$\pi^\star$. That row equals the corresponding row of some $p^{(\pi^\star)}\in\DLCN$; by
Definition~\ref{def:realizable} again (applied to $p^{(\pi^\star)}$, whose table has this same row)
$g_s(p^{(\pi^\star)})=g_s(q)$, and $p^{(\pi^\star)}\in\DLCN$ satisfies $s$. Hence
$g_s(q)\in[\ell_s,u_s]$.

If $s$ has form (R1), then $\pa_v=\varnothing$, so the family has the single (empty) parent
configuration and its table is the whole marginal $P(\ch_v)=P_{v,\varnothing}$. The same argument
applies verbatim with $\pi^\star$ the empty configuration.

In both cases the key point is that realizability makes $g_s$ a function of \emph{one row}, and a
row certified by (A) to be realized inside $\DLCN$ carries its constraint value with it --- so the
free, independent choice of the \emph{other} rows and families cannot disturb it. Therefore every
sentence holds at $q$, giving $q\in\DLCN$.
\end{proof}

\begin{remark}[Why there is no biconditional over ``every atom'']\label{rem:nobiconditional}
The companion note \texttt{strong\_extension\_\allowbreak exactness.tex} states its main theorem
as an \emph{iff}: $\DSE=\DLCN$ --- equivalently, Credal~VE is exact for every atom --- iff the primal
graph is singly connected and every sentence is family-realizable. Two problems make that
formulation untenable, and we therefore state Propositions~\ref{prop:incl}--\ref{prop:suff} plus
witness examples instead.

\emph{First, set-strictness does not imply marginal looseness.} On \texttt{examples/d4\_biting.lcn}
the cross-family sentence $P(B\wedge C)\le0.05$ is non-realizable and the free product genuinely
violates it (reaching $P(B\wedge C)\approx0.92$), so $\DSE\supsetneq\DLCN$. Yet all three
unconditional marginals are \emph{exact}: measured, Credal~VE returns
$P(A)=[0.40,0.60]$, $P(B)=[0.04,0.72]$, $P(C)=[0.08,0.80]$, matching
\texttt{ExactInference}(\texttt{global}) and \texttt{CredalJT} endpoint for endpoint. The extra
joints in $\DSE\setminus\DLCN$ simply do not attain any singleton-marginal extremum. So
``$\DSE=\DLCN$'' is strictly stronger than ``every atom is exact'', and the two cannot be equivalent.

\emph{Second, single-connectedness is not necessary} --- see Example~\ref{ex:loopyexact}.

What \emph{is} true, and is proved above, is sufficiency; the necessity of hypothesis~(B) is
established by the witnesses of Examples~\ref{ex:collider} and~\ref{ex:2node}, which show that
dropping realizability makes specific marginals strictly loose. That is a complete and honest account
of the boundary.
\end{remark}

\begin{example}[A loopy LCN on which Credal~VE is exact]\label{ex:loopyexact}
Single-connectedness is not necessary for exactness, and the topology table of earlier versions of
this note --- which marked Credal~VE ``outer only'' on \texttt{ktree-fr} \emph{because} that class is
loopy --- was wrong on this point. Consider
\texttt{examples/ktree\_fr\_k2\_n4\_example.lcn} ($n{=}4$, $k{=}2$):
\begin{center}\small
\begin{tabular}{@{}ll@{}}
$s_0$: $0.084382\le P(x_1)\le0.684382$ &
$s_2$: $0.092785\le P(x_3\mid x_1\wedge x_2)\le0.692785$\\
$s_1$: $0.01\le P(\lnot x_2\mid\lnot x_1)\le0.572656$ &
$s_3$: $0.068242\le P(x_0\mid x_1\wedge x_3)\le0.668242$
\end{tabular}
\end{center}
The seed $3$-clique is $\{x_1,x_2,x_3\}$ and $x_0$ attaches to the $2$-clique $\{x_1,x_3\}$, so the
moralized graph contains the triangle $x_1x_2x_3$: the instance is \textbf{loopy}, decidedly not
singly connected. It is also fully realizable (each $\psi$ is a full conjunction over its family's
parents; $s_0$ is a root marginal). Measured, with all eight SCIP solves \texttt{confirmed} at gap
$0$:
\begin{center}\small
\begin{tabular}{@{}lcccc@{}}
\toprule
engine & $P(x_0)$ & $P(x_1)$ & $P(x_2)$ & $P(x_3)$ \\
\midrule
\texttt{ExactInference} (global) & $[0,1]$ & $[0.0844,0.6844]$ & $[0.1349,0.9968]$ & $[0,1]$ \\
\texttt{CredalJT} (D5)           & $[0,1]$ & $[0.0844,0.6844]$ & $[0.1349,0.9968]$ & $[0,1]$ \\
\textbf{Credal~VE}               & $[0,1]$ & $[0.0844,0.6844]$ & $[0.1349,0.9968]$ & $[0,1]$ \\
\bottomrule
\end{tabular}
\end{center}
Credal~VE is \emph{exact on every atom} of a loopy LCN. Moreover \texttt{method="linear-tight"} (D1,
a pure Gap-A closer) returns identical numbers, so there is no residual Gap~A either. This is
consistent with Proposition~\ref{prop:suff}, whose hypotheses (A) and (B) hold here and which never
mentions the topology; and it refutes the inference ``loopy $\Rightarrow$ Gap~B''. The correct
statement is that loopiness alone does not create Gap~B: a non-realizable sentence, or an open
Gap~A, does.
\end{example}

\begin{remark}[What loopiness \emph{does} cost]\label{rem:loopycost}
Loopiness is not free --- but its price is paid in \emph{cost}, not accuracy. On the five-atom
\texttt{ktree\_fr\_k3\_n5\_1.lcn} of Example~\ref{ex:fr5}, the credal network builds essentially
instantly and is tiny ($8\cdot8\cdot4\cdot1\cdot2=512$ vertex combinations across the five families),
yet Credal~VE did not terminate within $200$\,s: it stalls inside \texttt{Potential.prune} during
elimination. \texttt{CredalJT} solves the same instance in about a second. So on bounded-treewidth
loopy classes the practical recommendation --- prefer \texttt{CredalJT} --- stands, but the reason is
tractability, not looseness.
\end{remark}

\section{The credal-network engines}\label{sec:cnfamily}

All engines in this section consume the same compiled \texttt{CredalNetworkVertices}; they differ only
in how they propagate the enumerated extreme points. Throughout, \texttt{method="linear"} is assumed
(Remark~\ref{rem:d1convex}).

\subsection{Two optimization lemmas}\label{sec:optlemmas}
Both results below rest on the same geometric fact, which we state carefully because the form usually
quoted (``a multilinear objective attains its extremum at a product of vertices'') is not the form
that is needed: the read-off is a \emph{ratio}, not a multilinear function.

\begin{lemma}[Block-affine vertex attainment]\label{lem:bauer}
Let $K=\prod_{j=1}^m K_j$ with each $K_j\subseteq\mathbb R^{d_j}$ compact and convex, and let
$F:K\to\mathbb R$ be continuous and \emph{affine in each block separately} (fixing all blocks but
$j$, the restriction is affine on $K_j$). Then $F$ attains its maximum and its minimum on
$\prod_j\ext(K_j)$.
\end{lemma}

\begin{proof}
$K$ is compact and $F$ continuous, so a maximizer $t^\star=(t_1^\star,\dots,t_m^\star)$ exists. We
show it can be moved to a product of extreme points without decreasing $F$. Fix $j$ and freeze the
other blocks at $t^\star$. The map $t_j\mapsto F(\dots,t_j,\dots)$ is affine on the compact convex
$K_j$, hence in particular convex; by the Bauer maximum principle a convex continuous function on a
compact convex set attains its maximum at an extreme point, so there is
$\hat t_j\in\ext(K_j)$ with $F(\dots,\hat t_j,\dots)\ge F(t^\star)$. Replace $t_j^\star$ by
$\hat t_j$; the value has not decreased. Iterating over $j=1,\dots,m$ (each step freezes the blocks
already replaced at extreme points, and affinity in block $j$ is unaffected by the values of the
others) yields a point of $\prod_j\ext(K_j)$ with value $\ge F(t^\star)$, hence equal to the maximum.
For the minimum apply the same argument to $-F$, which is also block-affine.
\end{proof}

\begin{lemma}[Block-fractional vertex attainment]\label{lem:fractional}
Let $K=\prod_{j=1}^m K_j$ with each $K_j$ compact convex, and let $N,D:K\to\mathbb R$ be continuous
and affine in each block separately, with $D(t)>0$ for all $t\in K$. Then $F=N/D$ attains its maximum
and its minimum on $\prod_j\ext(K_j)$.
\end{lemma}

\begin{proof}
$F$ is continuous on the compact $K$ (as $D>0$), so extrema exist. Fix $j$, freeze the other blocks,
and consider $\phi(t_j)=N(t_j)/D(t_j)$ on $K_j$, with $N,D$ affine in $t_j$ and $D>0$. Take any
segment $t_j(\theta)=(1-\theta)a+\theta b$, $\theta\in[0,1]$, in $K_j$. Then
$\theta\mapsto N(t_j(\theta))$ and $\theta\mapsto D(t_j(\theta))$ are affine in $\theta$, so
$\phi(t_j(\theta))=\frac{n_0+n_1\theta}{d_0+d_1\theta}$ with $d_0+d_1\theta>0$ on $[0,1]$. Its
derivative is $\frac{n_1d_0-n_0d_1}{(d_0+d_1\theta)^2}$, whose sign is constant on $[0,1]$: $\phi$ is
therefore \emph{monotone} along every segment. A function monotone along all segments of a convex set
attains both its maximum and its minimum at extreme points --- for the maximum: given any
$t_j\in K_j$, write it (Krein--Milman, finite dimension: Minkowski's theorem) as a convex combination
of extreme points and note that repeatedly moving along a segment towards whichever endpoint does not
decrease $\phi$ terminates at an extreme point with value $\ge\phi(t_j)$; symmetrically for the
minimum. So there is $\hat t_j\in\ext(K_j)$ not decreasing $F$. Iterate over the blocks as in
Lemma~\ref{lem:bauer}.
\end{proof}

\begin{remark}[Why the fractional form is the one needed]\label{rem:whyfrac}
Credal~VE's read-off (\texttt{cve.py}, Step~6) normalizes each surviving function,
$f\mapsto f/\sum_a f[a]$, and reports the componentwise extrema of the normalized vectors. So even
for an \emph{unconditional} marginal the reported quantity is a ratio of two block-affine functions,
and Lemma~\ref{lem:fractional} --- not Lemma~\ref{lem:bauer} --- is what applies. For a
\emph{conditional} query $P(a\mid e)=P(a,e)/P(e)$ the same lemma applies with $N,D$ the two
block-affine marginals. In both cases the hypothesis $D>0$ is essential and is \emph{not} automatic:
the implementation skips offending functions outright (\texttt{if total <= 0: continue}). We
therefore carry $\min_{P\in\DSE}P(e)>0$ as an explicit hypothesis.
\end{remark}

\subsection{Credal~VE}\label{sec:cve}
Credal~VE (\texttt{cn/cve.py}) computes each atom's bounds by a single-query bucket elimination, each
time recomputing an elimination order in which the target is eliminated \emph{last}. Potentials are
\emph{sets} of functions; \texttt{combine} takes elementwise products pairwise, \texttt{marginalize}
sums out a variable, and \texttt{Potential.prune} discards \emph{dominated} functions ($f$ is
dominated if some $g$ in the same potential satisfies $g\ge f$ componentwise, with strict inequality
somewhere).

\begin{lemma}[Prune monotonicity]\label{lem:prunemono}
Let $\Phi$ be any composition of \texttt{combine} operations (with fixed non-negative co-factors) and
\texttt{marginalize} operations, mapping functions on a scope $S$ to functions on a scope $S'$. If
$g\ge f$ pointwise on $S$ then $\Phi(g)\ge\Phi(f)$ pointwise on $S'$. Consequently, for every state
$s'$ of $S'$,
\[
  \max_{f\in\mathcal P}\Phi(f)[s']=\max_{f\in\mathrm{prune}(\mathcal P)}\Phi(f)[s'].
\]
\end{lemma}

\begin{proof}
It suffices to check the two primitive operations, since $\Phi$ is their composition. For
\texttt{combine} with a fixed $h\ge0$: $g\ge f$ implies $g\cdot h\ge f\cdot h$ pointwise, as
multiplying an inequality between non-negative numbers by a non-negative number preserves it. For
\texttt{marginalize} over a variable $x$: $g\ge f$ pointwise implies
$\sum_x g\ge\sum_x f$ pointwise in the remaining variables, as a sum of termwise inequalities.
Composing preserves the property.

For the second claim: \texttt{prune} removes $f$ only when some $g$ in the same potential dominates
it, and by the above $\Phi(g)[s']\ge\Phi(f)[s']$, so the removed $f$ cannot be the unique argmax at
any state $s'$. Since \texttt{prune} never removes all functions, and domination is transitive, the
maximum over the pruned set is unchanged.
\end{proof}

\begin{theorem}[Credal~VE is exact over the strong extension]\label{thm:cve}
Let $\mathcal L$ be a chain-graph LCN with \texttt{method="linear"} local sets, and run Credal~VE
with \texttt{coupling="off"} and exact pruning ($\epsilon{=}0$) on an \emph{unconditional} marginal
query $P(a)$. Then Credal~VE returns exactly the strong-extension interval,
\[
  [\lo q_{\mathrm{VE}}(a),\hi q_{\mathrm{VE}}(a)]
  =\Bigl[\min_{P\in\DSE}P(a),\ \max_{P\in\DSE}P(a)\Bigr],
\]
and hence, by Proposition~\ref{prop:incl}, a valid outer bound w.r.t.\ $\DLCN$; under the hypotheses
of Proposition~\ref{prop:suff} it is exact w.r.t.\ $\DLCN$.
\end{theorem}

\begin{proof}
Write $\mathcal V=\prod_v\prod_\pi\ext(K_{v,\pi})$ for the set of vertex selections, so that each
$\tau\in\mathcal V$ determines a product joint $P_\tau\in\DSE$.

\emph{Step 1: every function formed is a vertex product, and all have total mass one.} By induction
over the elimination schedule. Initially each family's potential is the set of its enumerated local
extreme points, indexed by the rows of that family. \texttt{combine} of two potentials forms all
pairwise elementwise products, and \texttt{marginalize} sums a variable out; hence every function
appearing at any stage is, for some $\tau\in\mathcal V$, the marginal of $P_\tau$ on the current
scope --- this is exactly the standard bucket-elimination identity for the precise network $P_\tau$,
applied uniformly across $\tau$. In particular, since the query is eliminated \emph{last}, when the
schedule finishes the surviving potential has scope $\{a\}$ and its functions are precisely
$\{(P_\tau\restr\{a\})\}_{\tau\in\mathcal T}$ for the set $\mathcal T\subseteq\mathcal V$ of
selections that survived pruning. Each such function is a probability distribution on $\{a\}$: it has
total mass $\sum_a P_\tau(a)=1$, because $P_\tau$ is a joint probability distribution and no evidence
indicator was multiplied in (this is where ``unconditional'' is used).

\emph{Step 2: normalization is the identity, so pruning is harmless.} With no evidence, every final
function has total mass $1$ by Step~1, so the read-off $f\mapsto f/\sum_a f[a]$ leaves each function
unchanged. Hence the reported bounds are the componentwise extrema
$\min_{\tau\in\mathcal T}P_\tau(a)$ and $\max_{\tau\in\mathcal T}P_\tau(a)$ of the surviving
functions. By Lemma~\ref{lem:prunemono} applied to the composition $\Phi$ of the operations performed
after each pruning step, these componentwise extrema over $\mathcal T$ coincide with those over the
full $\mathcal V$:
$\min_{\tau\in\mathcal T}P_\tau(a)=\min_{\tau\in\mathcal V}P_\tau(a)$ and likewise for $\max$.

\emph{Step 3: the vertex extremum is the $\DSE$ extremum.} The map
$\tau\mapsto P_\tau(a)$ extends to $\prod_v\prod_\pi K_{v,\pi}\to\mathbb R$,
$(P_{v,\pi})\mapsto\bigl(\prod_v\prod_\pi P_{v,\pi}\bigr)(a)$, which is affine in each block
$P_{v,\pi}$ separately: expanding the product, each block's table appears to degree exactly one. By
Lemma~\ref{lem:bauer} the extrema of this map over the product of the (compact, convex) $K_{v,\pi}$
are attained at $\mathcal V$. Since $\DSE$ is precisely the image of that product under the map,
\[
  \min_{P\in\DSE}P(a)=\min_{\tau\in\mathcal V}P_\tau(a),\qquad
  \max_{P\in\DSE}P(a)=\max_{\tau\in\mathcal V}P_\tau(a).
\]
Chaining Steps 2 and 3 gives the claim. The final sentence follows from
Propositions~\ref{prop:incl} and~\ref{prop:suff}.
\end{proof}

\begin{remark}[Conditional queries: a genuine gap]\label{rem:cvecond}
Theorem~\ref{thm:cve} is deliberately restricted to unconditional marginals, and the restriction
cannot simply be lifted. With evidence, an indicator potential is multiplied in, so a final function
has total mass $P_\tau(e)$, which \emph{varies with $\tau$}. Normalization is then a
$\tau$-\emph{dependent} rescaling and Step~2 breaks: the read-off ratio is not monotone under
componentwise domination. Concretely, over states $(a{=}1,a{=}0)$ the function $g=(1,1)$ dominates
$f=(1,0)$, so \texttt{prune} discards $f$ --- yet $f$ normalizes to read-off $1.0$ and $g$ to $0.5$.
Pruning can therefore discard the function carrying the extremal conditional, yielding an
\emph{inner} (unsound) interval. This is the same mechanism that makes Credal~CTE inner
(Section~\ref{sec:cte}); the difference is that CTE additionally shares one order across all queries,
so it is exposed even without evidence.

\emph{Measured:} the hazard did not materialize on the instances tested here. Monkey-patching
\texttt{Potential.prune} to the identity leaves Credal~VE's output bit-identical on
\texttt{d4\_biting}, on the collider witness of Example~\ref{ex:collider}, and on the two-atom
witness of Example~\ref{ex:2node}. So the theoretical gap is real but not currently observed; we
state it rather than assume it away.
\end{remark}

\begin{example}[Credal~VE against three references]\label{ex:cveex}
On \texttt{examples/d4\_biting.lcn} (three atoms $A\to B$, $A\to C$, plus the cross-family sentence
$P(B\wedge C)\le0.05$), Credal~VE returns $P(A)=[0.40,0.60]$, $P(B)=[0.04,0.72]$,
$P(C)=[0.08,0.80]$, agreeing with \texttt{ExactInference}(\texttt{global}) and \texttt{CredalJT} to
four decimals. Here $\DSE\supsetneq\DLCN$ (the free product violates the cap) yet the marginals
coincide --- the phenomenon of Remark~\ref{rem:nobiconditional}. On the loopy
Example~\ref{ex:loopyexact} Credal~VE is likewise exact on all four atoms; on the non-realizable
Examples~\ref{ex:collider} and~\ref{ex:2node} it returns the strictly wider strong-extension
interval, as Theorem~\ref{thm:cve} predicts.
\end{example}

\subsection{Credal~CTE: not exact, and not even outer}\label{sec:cte}
Credal~CTE (\texttt{cn/ccte.py}) is the two-pass (collect $+$ distribute) cluster-tree engine: one
\emph{shared} elimination order produces all marginals at once, deliberately \emph{not} augmented per
query. That single design decision costs it both exactness and soundness.

\begin{proposition}[CTE can be strictly inner]\label{prop:cte}
With \texttt{coupling="off"} and exact pruning, Credal~CTE's interval can be strictly contained in the
$\DSE$ interval; in particular it is \emph{not} in general a valid outer bound w.r.t.\ $\DLCN$.
\end{proposition}

\begin{proof}
It suffices to exhibit one instance, since the claim is a non-containment. On the shipped
\texttt{d4\_biting} instance the measured outputs are
\begin{center}\small
\begin{tabular}{@{}lccc@{}}
\toprule
engine & $P(A)$ & $P(B)$ & $P(C)$\\
\midrule
\texttt{ExactInference} (global, certified) & $[0.40,0.60]$ & $\mathbf{[0.04,0.72]}$ & $[0.08,0.80]$\\
Credal~VE $=$ \texttt{CredalJT} $=$ Interval~BP $=$ ApproxLP & $[0.40,0.60]$ & $[0.04,0.72]$ & $[0.08,0.80]$\\
\textbf{Credal~CTE} & $[0.40,0.60]$ & $\mathbf{[0.05,0.65]}$ & $[0.08,0.80]$\\
\bottomrule
\end{tabular}
\end{center}
Since $[0.05,0.65]\subsetneq[0.04,0.72]$ and the latter is the exact $\DLCN$ (and here also $\DSE$)
interval for $P(B)$, CTE's interval is strictly inner and excludes feasible marginals.

The mechanism is the interaction of Lemma~\ref{lem:prunemono}'s hypothesis with the shared order.
Credal~VE eliminates the query last, so by Step~2 of Theorem~\ref{thm:cve} the normalization at
read-off is trivial and dominance pruning is harmless. CTE eliminates $B$ mid-schedule, at which point
the function carrying the extremal $P(B)$ is dominated \emph{at that intermediate scope} and is
pruned --- but domination at a marginalized scope does not imply domination of the eventual
normalized read-off, exactly as in Remark~\ref{rem:cvecond}. The extremal vertex combination is
therefore discarded before it can be read off.
\end{proof}

\begin{remark}[A correction to the companion notes]\label{rem:ctecorrection}
\texttt{tighter\_approximation.tex} lists Credal~CTE as ``exact ($\epsilon{=}0$); outer if
$\epsilon>0$'', while \texttt{ccte\_exactness.tex} reports it as strictly inner. The measurement in
Proposition~\ref{prop:cte} settles the disagreement in favour of \texttt{ccte\_exactness.tex}. CTE
\emph{is} exact on a clean singly-connected, cross-family-sentence-free LCN (there the shared schedule
is a genuine clique tree and the surviving dominant function also dominates for every downstream
read-off), but that hypothesis must be stated: it is not exact at $\epsilon{=}0$ in general.
\end{remark}

\subsection{Interval~BP}\label{sec:ibp}
Interval~BP (\texttt{cn/ibp.py}) is interval message passing over the credal network's factor graph.
With \texttt{method="interval"} it propagates per-state intervals; with \texttt{method="variational"}
it runs mean-field over a finite sample of vertex combinations and is \emph{inner} by construction
(it evaluates a subset of $\DSE$).

\begin{theorem}[Interval~BP is outer, and exact on trees]\label{thm:ibp}
With \texttt{method="interval"}:
\begin{enumerate}
  \item[\textup{(i)}] For every $P\in\DSE$, every edge and every iteration, $P$'s marginal on the
    message's variable lies in the message interval. Hence the returned interval is a valid
    \emph{outer} bound w.r.t.\ $\DSE$, on any topology.
  \item[\textup{(ii)}] If the factor graph is a tree, the converged intervals equal the $\DSE$
    interval for every atom; combined with Proposition~\ref{prop:suff} they are then exact w.r.t.\
    $\DLCN$.
\end{enumerate}
\end{theorem}

\begin{proof}
(i) Induction on the iteration count $t$, simultaneously over all edges. At $t=0$ every message is
$[0,1]$, which contains every marginal. Assume the claim at $t$ and fix $P=\prod_v\prod_\pi P_{v,\pi}
\in\DSE$. A \emph{variable-to-factor} message is the intersection of the other incoming
factor-to-variable intervals; by the inductive hypothesis each of those contains $P$'s marginal of
that variable, so their intersection does too (an intersection of sets each containing a common point
contains it). A \emph{factor-to-child} update ranges over all local extreme points of that factor's
credal set and over all corner distributions of the incoming messages, and returns the elementwise
min/max of the resulting marginal. Since $P$'s own local table at this factor is a convex combination
of that factor's extreme points, and $P$'s incoming marginals lie in the incoming intervals by
induction, $P$'s marginal of the child is a convex combination of values that the update ranged over,
hence lies between the min and max it reports. A conditional factor sends the vacuous $[0,1]$ upward,
which excludes nothing. So the claim holds at $t+1$. Taking the final intersection over incident
factors preserves containment, giving the outer bound.

(ii) Suppose the factor graph is a tree. We show the converged interval is also \emph{attained},
which with (i) gives equality. Root the tree at the query atom $a$ and induct on subtree depth. For a
leaf factor the update optimizes over that factor's extreme points with no incoming constraint, so
both endpoints are attained by an actual vertex selection. For an internal factor $f$ with children
subtrees $T_1,\dots,T_k$: because the graph is a tree, the subtrees are vertex-disjoint apart from $f$
and share no variable, so the vertex selections in distinct $T_i$ are \emph{independent} degrees of
freedom of $\DSE$ --- no selection is constrained twice. By induction each incoming message endpoint
is attained by some selection within its subtree; combining those selections (legitimate, by
independence) with the extreme point of $f$'s own set that the update chose yields a single
$\tau\in\mathcal V$ attaining the reported endpoint. Since the message endpoints are attained rather
than merely containing, no slack is introduced, and at the root the reported interval equals
$\{P_\tau(a)\}$'s extremes, which by Step~3 of Theorem~\ref{thm:cve} is the $\DSE$ interval. On a
loopy graph the independence step fails --- a selection reachable through two paths would be
constrained twice while the interval algebra treats the two messages as independent --- which is
precisely the source of loop slack.
\end{proof}

\begin{remark}[A residual architectural limitation, measured]\label{rem:ibpevid}
Theorem~\ref{thm:ibp}(ii) concerns unconditional marginals. Interval~BP's intersection message algebra
performs no Bayesian \emph{upward} update, so evidence on a child does not shift a parent's posterior.
On the two-atom chain of Example~\ref{ex:2node} (without \texttt{M1}), querying $P(x_0\mid x_1{=}1)$:
\begin{center}\small
\begin{tabular}{@{}lc@{}}
\toprule
engine & $P(x_0\mid x_1{=}1)$\\
\midrule
Interval~BP (\texttt{interval}) & $[0.30,0.50]$ --- the \emph{unchanged prior}\\
Credal~VE & $[0.1765,0.3333]$\\
\texttt{CredalJT} (D5) & $[0.0968,0.4000]$\\
\bottomrule
\end{tabular}
\end{center}
Interval~BP simply returns $P(x_0)$: the evidence has no effect. For conditional queries with
evidence on descendants, use Credal~VE or \texttt{CredalJT}. (That the latter two also disagree here
is expected --- Credal~VE bounds $\DSE$ conditionally, subject to Remark~\ref{rem:cvecond}, while
\texttt{CredalJT} targets $\DLCN$.)
\end{remark}

\begin{remark}[Shipped status: two bugs, both fixed]\label{rem:ibpbugs}
Interval~BP previously carried two message-passing bugs that made the interval mode \emph{inner} and
unsound, contradicting Theorem~\ref{thm:ibp}(i): a spurious normalized-likelihood upward message
(which collapsed an unconstrained root to $[0.5,0.5]$) and an invalid component-wise corner
enumeration. Both are corrected. Verified: on the two-atom chain the engine now returns
$P(x_0)=[0.30,0.50]$ (was $[0.5,0.5]$) and on \texttt{d4\_biting} $P(B)=[0.04,0.72]$ (was
$[0.05,0.65]$).
\end{remark}

\subsection{ApproxLP}\label{sec:approxlp}
ApproxLP (\texttt{cn/approxlp.py}) solves the multilinear program by coordinate descent: initialize
every local table to the \emph{center} of its credal set (the mean of the enumerated extreme points);
then sweep, and for each (family, parent-configuration) block hold the others fixed, enumerate that
block's extreme points, evaluate the objective exactly by variable elimination on the resulting
\emph{precise} Bayesian network, and greedily commit the best; repeat to convergence.

\begin{theorem}[ApproxLP is inner w.r.t.\ the strong extension]\label{thm:alp}
Run ApproxLP with \texttt{coupling="off"} on a query whose denominator is positive throughout (for an
unconditional marginal this is automatic). Then
\[
  \min_{P\in\DSE}q(P)\;\le\;\lo q_{\mathrm{ALP}}\;\le\;\hi q_{\mathrm{ALP}}\;\le\;\max_{P\in\DSE}q(P),
\]
i.e.\ ApproxLP's interval is contained in the $\DSE$ interval. It is therefore not in general a valid
outer bound w.r.t.\ $\DSE$.
\end{theorem}

\begin{proof}
\emph{Step 1: every state visited is a point of $\DSE$.} By induction over the sweep. The initial
state assigns each block the center of $K_{v,\pi}$, i.e.\ the arithmetic mean of its extreme points;
since $K_{v,\pi}$ is convex and its extreme points belong to it, the mean lies in $K_{v,\pi}$. For the
inductive step, a commit replaces one block's value by an extreme point of that same $K_{v,\pi}$,
which lies in $K_{v,\pi}$ by definition; all other blocks are untouched. Hence after every commit the
state is a tuple $(P_{v,\pi})$ with $P_{v,\pi}\in K_{v,\pi}$ for all $(v,\pi)$ --- note that this is
exactly membership of the assembled product in $\DSE$ as defined by \eqref{eq:se}, because
\eqref{eq:se} allows each block to be chosen independently (Remark~\ref{rem:rowwise}). Mixed states,
with some blocks at centers and others at vertices, are therefore admissible members of $\DSE$; this
is the point at which the row-wise reading of \eqref{eq:se} is load-bearing.

\emph{Step 2: the objective is evaluated exactly at each visited state.} At a visited state the
network is a \emph{precise} Bayesian network $P_\tau$ (every block has a single table), and the
evaluation runs standard variable elimination on the factorization \eqref{eq:factor}. Variable
elimination computes marginals of a precise factorized distribution exactly (it is the distributive
law applied to the sum over eliminated variables), and for a conditional query it forms the ratio of
two such marginals, which is exact provided the denominator $P_\tau(e)>0$ --- the hypothesis. Hence
every value ApproxLP records equals $q(P)$ for that specific $P\in\DSE$.

\emph{Step 3: conclusion.} By Steps 1--2 the set of values ApproxLP ever records is a subset of
$\{q(P):P\in\DSE\}$. Its reported lower bound is the value at the terminal state of the
\texttt{sense="min"} descent and its reported upper bound that of the \texttt{sense="max"} descent;
both are elements of that value set. A minimum over a subset is $\ge$ the minimum over the whole set,
and a maximum over a subset is $\le$ the maximum, giving
$\min_{\DSE}q\le\lo q_{\mathrm{ALP}}$ and $\hi q_{\mathrm{ALP}}\le\max_{\DSE}q$. Finally
$\lo q_{\mathrm{ALP}}\le\hi q_{\mathrm{ALP}}$: both descents start from the \emph{same} center state
$\tau_0$, and each is monotone in its own direction (a commit is made only if it improves the
objective), so $\lo q_{\mathrm{ALP}}\le q(\tau_0)\le\hi q_{\mathrm{ALP}}$.
\end{proof}

\begin{proposition}[Tightness is an assumption about the landscape, not a topological condition]
\label{prop:alptight}
Suppose (a) $\DSE=\DLCN$ (e.g.\ under the hypotheses of Proposition~\ref{prop:suff}), and (b) for
each sense the coordinate descent terminates at a \emph{global} optimum of $q$ over $\DSE$. Then
ApproxLP's interval equals the exact $\DLCN$ interval.
\end{proposition}

\begin{proof}
Immediate from Theorem~\ref{thm:alp}: by (b) the two terminal values are the global $\DSE$ extrema, so
the chain of inequalities in Theorem~\ref{thm:alp} collapses to equalities, and by (a) the $\DSE$
interval is the $\DLCN$ interval.
\end{proof}

\begin{remark}[Honesty about hypothesis (b)]\label{rem:alphonest}
Hypothesis (b) is an assumption about the optimization landscape, not a checkable structural
condition, and we state it as such. The companion note's version of this result assumes instead that
``the per-family linearizations expose their global optimum to a greedy sweep from the center'', and
then proves tightness by invoking that assumption --- which is the conclusion. No structural
sufficient condition for (b) is known; the multilinear program is nonconvex, so a sweep can stall
where no single-block move improves but a coordinated multi-block jump would. The practical
substitute is the bracket of Corollary~\ref{cor:bracket}.
\end{remark}

\begin{corollary}[Two-sided bracket and stall detection]\label{cor:bracket}
Run ApproxLP and Credal~VE on the \emph{same} enumerated vertices for an unconditional marginal. Then
$[\lo q_{\mathrm{ALP}},\hi q_{\mathrm{ALP}}]\subseteq[\lo q_{\mathrm{SE}},\hi q_{\mathrm{SE}}]
=[\lo q_{\mathrm{VE}},\hi q_{\mathrm{VE}}]$. If the two intervals coincide, the common value is
certified to be the exact $\DSE$ bound; if they differ, the difference is exactly the
coordinate-descent stall.
\end{corollary}

\begin{proof}
The inclusion is Theorem~\ref{thm:alp} and the equality is Theorem~\ref{thm:cve}. If the endpoints
coincide then ApproxLP attained the $\DSE$ extremum, certifying it; otherwise
Theorem~\ref{thm:alp}'s inequality is strict, which by Step~3 of its proof happens only when the
terminal state is not a global optimum.
\end{proof}

\begin{example}[ApproxLP is inner w.r.t.\ $\DSE$ yet \emph{outer} w.r.t.\ $\DLCN$]\label{ex:alpouter}
This is why Definition~\ref{def:char} insists on naming the reference set. On \texttt{d4\_biting}
ApproxLP returns $P(B)=[0.04,0.72]$ --- coinciding with Credal~VE, so by
Corollary~\ref{cor:bracket} the $\DSE$ bound is certified and the inner gap is zero. But on the
two-atom witness of Example~\ref{ex:2node} \emph{with} \texttt{M1}, ApproxLP returns
$P(x_1)=[0.40,0.68]$ while the exact $\DLCN$ bound is $[0.45,0.55]$: here ApproxLP strictly
\emph{contains} the truth, i.e.\ it is an \emph{outer} bound w.r.t.\ $\DLCN$. Both facts are
consistent --- ApproxLP is inner w.r.t.\ $\DSE$, and $\DLCN\subsetneq\DSE$ here, with the Gap~B
looseness dominating the (zero) descent slack. The unqualified assertion ``ApproxLP is not a valid
outer bound'' is therefore false as stated; the correct statement is that it carries no outer
guarantee w.r.t.\ $\DSE$, and its position relative to $\DLCN$ depends on which effect dominates.
\end{example}

\subsection{Credal~IJGP}\label{sec:ijgp}
Credal~IJGP (\texttt{cn/ijgp.py}) runs iterative join-graph propagation with a mini-bucket
\texttt{i\_bound} capping cluster size.

\begin{proposition}[IJGP at large $i$-bound]\label{prop:ijgp}
If \texttt{i\_bound} $\ge$ the treewidth of the credal network's moral graph, no mini-bucket
partitioning is forced, the join graph is a junction tree, and Credal~IJGP coincides with exact
elimination over the enumerated vertices --- hence is exact over $\DSE$. For smaller
\texttt{i\_bound} it is an outer-style approximation whose slack the $i$-bound trades against a cost
of $2^{i}$ per cluster.
\end{proposition}

\begin{proof}
At \texttt{i\_bound} $\ge$ treewidth every bucket fits within the cap, so the mini-bucket splitting
step is vacuous and the join graph produced is the junction tree of the elimination order, which has
the running-intersection property and no cycle. Message passing on a cycle-free join graph is exact
elimination (the same argument as Theorem~\ref{thm:ibp}(ii), one level up at cluster rather than
variable granularity: disjoint subtrees supply independent vertex selections, so each message endpoint
is attained). By Step~3 of Theorem~\ref{thm:cve} the vertex extremum is the $\DSE$ extremum. When
\texttt{i\_bound} is below the treewidth, a bucket is split into mini-buckets whose messages are
propagated independently, which duplicates the influence of shared selections and introduces the
usual loopy slack.
\end{proof}

\begin{example}[IJGP on \texttt{d4\_biting}]\label{ex:ijgp}
Measured: $P(A)=[0.40,0.60]$, $P(B)=[0.04,0.72]$, $P(C)=[0.08,0.80]$ at both \texttt{i\_bound}$=1$
and \texttt{i\_bound}$=3$ --- exact, matching the certified oracle. The instance's moral graph is the
single clique $\{A,B,C\}$ of treewidth $2$, and even the $i$-bound-$1$ run happens not to lose
anything here; on larger loopy instances the $i$-bound is the tightness dial of
Proposition~\ref{prop:ijgp}.
\end{example}

\section{CredalJT: exact at treewidth cost}\label{sec:d5}

\texttt{CredalJT} (\texttt{cn/junction\_nlp.py}, scheme D5) does not compile to a credal network at
all: it builds one junction tree from the chain-graph moral graph, gives each cluster a variable
vector over the joint of its atoms, ties adjacent clusters by separator-consistency equalities, hosts
every sentence and LMC equality on a cluster containing its atoms, and optimizes the read-off with
SCIP. Its cost is $\sum_C 2^{|C|}=O(n\,2^{w+1})$ --- exponential in the \emph{treewidth} $w$, not in
$n$.

\begin{definition}[Cluster-feasible set, hosting, separator-entailment, read-off]\label{def:d5}
Let $T$ be a clique tree over clusters $\mathcal C$, built by the standard
moralize~$\to$~triangulate~$\to$~clique-tree pipeline, and let $G^\ast$ be the chordal graph whose
maximal cliques are the clusters (two atoms adjacent iff they co-occur in a cluster).
\begin{itemize}
  \item The \emph{cluster-feasible set} is
  \begin{align*}
    \mathcal G=\bigl\{(q_C)_{C\in\mathcal C}\;:\;
      &q_C\in\Delta(2^{|C|})\ \text{for each }C;\\
      &q_C\restr S=q_D\restr S\ \text{for every edge }(C,D)\in T,\ S=C\cap D;\\
      &\text{every hosted constraint holds on its host cluster}\bigr\}.
  \end{align*}
  \item A constraint $c$ is \emph{hosted} if some cluster $C_h\supseteq\mathrm{scope}(c)$ exists; its
    row is then imposed on $q_{C_h}$'s marginal.
  \item An LMC assertion $(X\indep Y\mid S)$ is \emph{separator-entailed} if $S$ separates $X$ from
    $Y$ in $G^\ast$.
  \item For a query atom $a$ and evidence $e$ lying together in a cluster $C_a$, the \emph{read-off}
    is
    $\readoff_a(q)=\sum_{i:a=1}q_{C_a}[i]$ (no evidence), or
    $\readoff_a(q)=\bigl(\sum_{i:a=1,e}q_{C_a}[i]\bigr)/\bigl(\sum_{i:e}q_{C_a}[i]\bigr)$ (with
    evidence).
  \item $\pi(p)=(p\restr C)_{C\in\mathcal C}$ is the \emph{projection} of a joint onto the clusters.
\end{itemize}
\end{definition}

The reconstruction direction needs the following lemma, whose usual statement contains two defects we
must repair: it is normally asserted with a \emph{uniqueness} claim that is false in the presence of
zero separator cells, and the division in the product formula is left unaddressed.

\begin{lemma}[RIP reconstruction, with zeros]\label{lem:rip}
Let $T$ be a clique tree with the running-intersection property over clusters $\mathcal C$, and let
$(q_C)\in\prod_C\Delta(2^{|C|})$ be separator-consistent: $q_C\restr S=q_D\restr S$ for every edge
$(C,D)$ with $S=C\cap D$. Then there exists a joint distribution $p\in\Delta(2^{\mathbf V})$ with
$p\restr C=q_C$ for every $C\in\mathcal C$, and $p$ factorizes over the cliques of $G^\ast$. Such $p$
is unique if every separator marginal is strictly positive, but \emph{not} in general.
\end{lemma}

\begin{proof}
Induction on $|\mathcal C|$. If $|\mathcal C|=1$ take $p=q_C$.

Let $|\mathcal C|\ge2$ and pick a leaf cluster $C$ of $T$ with neighbour $D$ and separator
$S=C\cap D$. Let $T'=T-C$; by the running-intersection property $C\setminus S$ meets no cluster of
$T'$, and $(q_{C'})_{C'\in T'}$ is still separator-consistent, so the inductive hypothesis gives a
joint $p'$ over $\mathbf V'=\bigcup_{C'\in T'}C'$ with $p'\restr C'=q_{C'}$ for all $C'\in T'$. Note
$p'\restr S=(p'\restr D)\restr S=q_D\restr S=q_C\restr S=:q_S$, using separator consistency.

Define a conditional kernel on the fresh atoms $C\setminus S$: for each assignment $s$ of $S$,
\[
  \kappa(\,\cdot\mid s)=
  \begin{cases}
    q_C(\,\cdot\,,s)/q_S(s), & q_S(s)>0,\\
    \delta_{c_0}, & q_S(s)=0,
  \end{cases}
\]
where $c_0$ is any fixed assignment of $C\setminus S$. This is well defined: when $q_S(s)>0$ the
right-hand side is non-negative and sums over $C\setminus S$ to $q_S(s)/q_S(s)=1$, so it is a
distribution. Now set $p(x)=p'(x_{\mathbf V'})\cdot\kappa(x_{C\setminus S}\mid x_S)$.

$p$ is a distribution: it is non-negative, and summing over $x_{C\setminus S}$ first gives $p'$, which
sums to $1$.

\emph{The arbitrary choice is irrelevant.} On $\{x:q_S(x_S)=0\}$ we have $p'(x_{\mathbf V'})=0$,
because $p'\restr S=q_S$ assigns zero mass to that event and $p'\ge0$. Hence $p$ vanishes there
regardless of $\kappa$, so $p$ does not depend on $c_0$. (This is also the precise sense in which the
familiar quotient formula $p=\prod_C q_C/\prod_{(C,D)}q_S$ is legitimate: a zero separator cell forces
every contributing cluster cell to zero, since $q_S=q_C\restr S$ is a sum of non-negative $q_C$
entries, so numerator and denominator vanish together and the term is assigned $0$.)

\emph{Cluster marginals.} For $C'\in T'$: summing $p$ over $x_{C\setminus S}$ gives $p'$, so
$p\restr C'=p'\restr C'=q_{C'}$. For $C$ itself: $p\restr C(c,s)=\bigl(\sum$ of $p$ over
$\mathbf V'\setminus S$ at fixed $x_S=s\bigr)\cdot\kappa(c\mid s)=p'\restr S(s)\,\kappa(c\mid s)
=q_S(s)\kappa(c\mid s)$, which equals $q_C(c,s)$ when $q_S(s)>0$ and equals $0=q_C(c,s)$ when
$q_S(s)=0$. So $p\restr C=q_C$.

\emph{Factorization.} Unwinding the induction, $p$ is a product of one non-negative factor per cluster
(namely $q_C$ or the kernel $\kappa$ derived from it, each a function of the atoms of that cluster
only), hence factorizes over the cliques of $G^\ast$.

\emph{Non-uniqueness.} If some separator cell $q_S(s)=0$, then $\kappa(\cdot\mid s)$ was chosen
arbitrarily; any other choice yields a joint with the same cluster marginals, since all such joints
vanish on that event. Hence $p$ is not unique. When all separator marginals are positive the kernels
are forced and the induction determines $p$ uniquely.
\end{proof}

\begin{theorem}[CredalJT exactness]\label{thm:d5}
Suppose \textup{(H1)} $T$ has the running-intersection property (true by construction);
\textup{(H2)} every LCN sentence is hosted; \textup{(H3)} every LMC assertion is either hosted or
separator-entailed in $G^\ast$; \textup{(H4)} the query atom and all evidence atoms lie together in
some cluster. Then for every such atom $a$,
\[
  \min_{q\in\mathcal G}\readoff_a(q)=\min_{p\in\DLCN}\readoff_a(\pi(p))=\lo P(a),
  \qquad
  \max_{q\in\mathcal G}\readoff_a(q)=\max_{p\in\DLCN}\readoff_a(\pi(p))=\hi P(a),
\]
and likewise for a conditional query $P(a\mid e)$ subject to Remark~\ref{rem:denfloor}. Moreover the
\emph{projection} half holds without \textup{(H2)}/\textup{(H3)}, so \texttt{CredalJT} is always a
valid \emph{outer} bound.
\end{theorem}

\begin{proof}
First note the read-off identity: for any joint $p$, $\readoff_a(\pi(p))=P_p(a)$ (resp.\
$P_p(a\mid e)$), because $\readoff_a$ sums the cells of $\pi(p)_{C_a}=p\restr C_a$ in which $a$ (and
$e$) hold, which by definition of a marginal equals $P_p(a)$ (resp.\ the ratio). So both sides
optimize the \emph{same} functional.

\emph{Projection $\DLCN\to\mathcal G$ (needs only H1, H4).} Let $p\in\DLCN$ and set
$q_C:=p\restr C$, i.e.\ $q=\pi(p)$. Each $q_C$ is a distribution (a marginal of one). Each separator
equality holds because both $q_C\restr S$ and $q_D\restr S$ equal $p\restr S$. Each hosted constraint
$c$ holds: by Lemma~\ref{lem:locality} its value depends on $p$ only through
$p\restr\mathrm{scope}(c)$, and since $\mathrm{scope}(c)\subseteq C_h$ we have
$q_{C_h}\restr\mathrm{scope}(c)=p\restr\mathrm{scope}(c)$, so $c$ holds on the host exactly because it
holds for $p$. Hence $q\in\mathcal G$, with $\readoff_a(q)=P_p(a)$. Therefore every value achievable
over $\DLCN$ is achievable over $\mathcal G$, giving
$\min_{\mathcal G}\readoff_a\le\min_{\DLCN}$ and $\max_{\mathcal G}\readoff_a\ge\max_{\DLCN}$: the
cluster interval \emph{contains} the true interval. This is the outer-bound claim, and it used neither
H2 nor H3.

\emph{Reconstruction $\mathcal G\to\DLCN$ (needs H2, H3).} Let $q\in\mathcal G$. By H1 and
Lemma~\ref{lem:rip} there is a joint $p$ with $p\restr C=q_C$ for every cluster, factorizing over the
cliques of $G^\ast$. We check $p\in\DLCN$.

Sub-scope consistency: if $A\subseteq C$ for some cluster $C$, then $p\restr A=(p\restr C)\restr A
=q_C\restr A$, since marginalizing in stages agrees with marginalizing at once.

\emph{Hosted constraints} (H2 and the hosted part of H3): let $c$ be hosted on $C_h$. By the previous
paragraph $p\restr\mathrm{scope}(c)=q_{C_h}\restr\mathrm{scope}(c)$, and $q$ satisfies $c$ on
$C_h$ by definition of $\mathcal G$; by Lemma~\ref{lem:locality} $c$ therefore holds for $p$.

\emph{Separator-entailed LMC} (the other part of H3): let $(X\indep Y\mid S)$ be separator-entailed,
so $S$ separates $X$ from $Y$ in $G^\ast$. Since $p$ factorizes over the cliques of $G^\ast$, it
satisfies the \emph{global Markov property} of $G^\ast$: factorization implies global Markov for any
undirected graph, and crucially this direction requires \emph{no} positivity assumption (Lauritzen,
\emph{Graphical Models}, Prop.~3.8; only the converse, global Markov $\Rightarrow$ factorization, needs
$p>0$ via Hammersley--Clifford). Global Markov states that if $S$ separates $A$ from $B$ in $G^\ast$
then $A\indep B\mid S$ under $p$; applying it to the connected components separated by $S$ and then
restricting to the subsets $X$ and $Y$ of those components (conditional independence is preserved
under taking subsets, by the decomposition property of conditional independence) gives
$X\indep Y\mid S$ under $p$, i.e.\ \eqref{eq:lmc} holds. Note this assertion was never imposed on
$\mathcal G$: it holds automatically for the reconstructed $p$.

Hence $p\in\DLCN$, and $\readoff_a(q)=\readoff_a(\pi(p))=P_p(a)$ because $\pi(p)=q$ on the host
cluster. So every value achievable over $\mathcal G$ is achievable over $\DLCN$, giving the reverse
inequalities. Combined with the projection half, the optima coincide. For a conditional query, by H4
the read-off is a functional of the single cluster $C_a$ and both maps preserve $q_{C_a}$, so the same
argument applies verbatim.
\end{proof}

\begin{remark}[Reading the proof]\label{rem:d5reading}
(a) The argument never mentions chains: exactness is a property of the \emph{junction tree}, via RIP,
hosting, separator-entailment and locality. (b) The projection half is unconditional, so
\texttt{CredalJT} is \emph{never unsound}: its interval always contains the truth, and H2/H3 decide
only tightness. (c) H3 is \emph{sufficient, not necessary}: on \texttt{examples/asia.lcn} and
\texttt{examples/polytree.lcn} some assertions are neither hosted nor separator-entailed, yet
\texttt{CredalJT} is empirically exact --- the unentailed assertions simply do not bind. Where H3
fails \emph{and} binds, the result is a valid but loose outer bound: on \texttt{examples/smokers.lcn}
the six large-scope LMCs of the undirected component $F_1F_2F_3$ are unhostable and unentailed
(within the width-$6$ cluster $\{F_1,F_2,F_3,S_1,S_2,S_3\}$ no separator divides $S_3$ from $F_1$),
and \texttt{CredalJT} returns $P(F_1)=[0,1]$.
\end{remark}

\begin{remark}[The denominator floor is a soundness hypothesis]\label{rem:denfloor}
For conditional queries the implementation linearizes the ratio with an auxiliary variable and, under
SCIP, adds a denominator floor $\mathbf 1_e^\top q_{C_a}\ge\varepsilon$ with
\texttt{\_DEN\_FLOOR}$=10^{-6}$ (\texttt{junction\_nlp.py:74,662}). This is an extra constraint on
$\mathcal G$ that $\DLCN$ does not have, so the projection half of Theorem~\ref{thm:d5} --- and with
it the outer-bound guarantee --- holds only for joints with $P(e)\ge\varepsilon$. If
$\min_{p\in\DLCN}P(e)<\varepsilon$ the returned conditional interval can be \emph{unsound}. This is a
soundness caveat, not merely a tightness one, and it is the exact analogue of Credal~VE's
\texttt{if total <= 0: continue} (Remark~\ref{rem:whyfrac}).
\end{remark}

\begin{corollary}[Exactness on \texttt{ktree-fr} at treewidth cost]\label{cor:ktree}
On every \texttt{ktree-fr} instance, H1--H4 hold, so \texttt{CredalJT} returns the exact $\DLCN$
marginal bounds at cost $O(n\,2^{k+1})$.
\end{corollary}

\begin{proof}
A $k$-tree is chordal, and the generator renders it as a DAG in construction order in which every atom
conditions on the \emph{full conjunction} of its $k$ clique-parents. Those parents are pairwise
adjacent in the $k$-tree, so moralization adds no edge and triangulation is a no-op: $G^\ast$ is the
$k$-tree itself and its maximal cliques --- the clusters --- are exactly the $(k{+}1)$-cliques, giving
H1 by construction. Every sentence is either a root marginal or scoped to a child with its $\le k$
parents, which co-occur in a cluster: H2. Each LMC assertion produced by the Local Markov Condition
has as its conditioning set a full clique-separator of the $k$-tree, which separates the two sides in
$G^\ast$, so each is hosted or separator-entailed: H3. Any single query atom lies in a
$(k{+}1)$-clique: H4. Theorem~\ref{thm:d5} applies, and the cost is $\sum_C2^{|C|}$ over $O(n)$
clusters each of size $k{+}1$.
\end{proof}

\begin{remark}[What the \texttt{ktree-fr} corollary does and does not claim]\label{rem:ktreehonest}
Corollary~\ref{cor:ktree} is a statement about \emph{cost}: exactness at $O(n2^{k+1})$ rather than
$2^n$. It should \emph{not} be read as ``the one loopy topology where a treewidth-bounded engine beats
the strong-extension engines on accuracy''. Example~\ref{ex:loopyexact} shows Credal~VE is also exact
on a loopy realizable $k$-tree; on the larger $k{=}3$ instance of Example~\ref{ex:fr5} Credal~VE fails
not by being loose but by not finishing (Remark~\ref{rem:loopycost}). The separation between
\texttt{CredalJT} and Credal~VE on this family is one of tractability.
\end{remark}

\begin{example}[Verified end to end]\label{ex:d5verified}
Example~\ref{ex:fr5} is the worked case: on \texttt{ktree\_fr\_k3\_n5\_1.lcn} the two maximal cliques
are $\{x_0,x_1,x_3,x_4\}$ and $\{x_1,x_2,x_3,x_4\}$ (treewidth $3$), the sole LMC assertion
$(x_0\indep x_2\mid x_1,x_3,x_4)$ has conditioning set equal to the separator between them, and
\texttt{CredalJT} reproduces the certified \texttt{ExactInference}(\texttt{global}) bounds on all five
atoms. Further checks: on \texttt{examples/chain.lcn} the clusters are the $7$ edge-pairs
(treewidth $1$: $28$ NLP variables against a $2^8=256$ joint); on \texttt{examples/polytree.lcn} the
collider cliques $\{x_0,x_1,x_2\}$ and $\{x_3,x_4,x_5\}$ give treewidth $2$; and on the loopy
Example~\ref{ex:loopyexact} \texttt{CredalJT} matches the certified oracle on every atom.
\end{example}

\section{ARIEL}\label{sec:ariel}

ARIEL (\texttt{lcn/inference/marginal/ariel.py}) is included here for reference: it targets $\DLCN$
directly on the primal-graph factor graph, without compiling a credal network, so it sits outside the
$\DSE$ family.

\subsection{The algorithm}
A \emph{factor node} groups all sentences sharing the \emph{same} atom scope and is joined to each
variable node in that scope (\texttt{utils/factor\_graph.py}). Two message kinds are passed to a
fixed point:
\begin{itemize}
  \item \emph{Variable-to-factor} $v\!\to\!f$: the intersection of the other incoming intervals,
    $\lo\ell_{v\to f}=\max_{f'\ne f}\lo\ell_{f'\to v}$ and
    $\hi u_{v\to f}=\min_{f'\ne f}\hi u_{f'\to v}$.
  \item \emph{Factor-to-variable} $f\!\to\!n$: minimize and maximize $P(n)$ over a \emph{local} NLP on
    the factor's own joint $\vec p\in\Delta(2^{|\mathrm{scope}(f)|})$, subject to (i) $f$'s sentences;
    (ii) the LMC equalities \eqref{eq:lmc} whose \emph{entire} scope fits inside
    $\mathrm{scope}(f)$; and (iii) for each other boundary variable $m$, the box constraint
    $\lo\ell_{m\to f}\le P(m)\le\hi u_{m\to f}$.
\end{itemize}
The answer for atom $v$ is the intersection $\lo P(v)=\max_f\lo\ell_{f\to v}$,
$\hi P(v)=\min_f\hi u_{f\to v}$.

\begin{lemma}[Message soundness]\label{lem:arielsound}
For every $P^\star\in\DLCN$, every edge and every iteration $t$, $P^\star$'s marginal of the message's
variable lies in the message interval. Consequently ARIEL's returned interval is a valid \emph{outer}
bound w.r.t.\ $\DLCN$ --- on \emph{any} topology, loopy included.
\end{lemma}

\begin{proof}
Induction on $t$, simultaneously over all edges. At $t=0$ every message is $[0,1]$.

Assume the claim at $t$. For a \emph{variable-to-factor} message, the update is an intersection of
intervals each of which contains $P^\star(v)$ by the inductive hypothesis, hence so does the
intersection.

For a \emph{factor-to-variable} message $f\!\to\!n$, consider the point
$\vec p^\star:=P^\star\restr\mathrm{scope}(f)$. We claim it is feasible for $f$'s local NLP. It lies in
the simplex, being a marginal of a distribution. It satisfies $f$'s sentences: each such sentence has
$\atoms(s)\subseteq\mathrm{scope}(f)$ and holds for $P^\star$ (as $P^\star\in\DLCN$), so by
Lemma~\ref{lem:locality} it holds for $\vec p^\star$. It satisfies the in-scope LMC equalities by the
same locality argument. And for each other boundary variable $m$, the box constraint
$\lo\ell_{m\to f}\le P(m)\le\hi u_{m\to f}$ is satisfied because
$\vec p^\star$'s marginal of $m$ is $P^\star(m)$, which lies in the incoming interval by the inductive
hypothesis. Hence $\vec p^\star$ is feasible, so the minimum of $P(n)$ over the local NLP is
$\le P^\star(n)$ and the maximum is $\ge P^\star(n)$: the new message contains $P^\star(n)$.

The final read-off intersects intervals all containing $P^\star(v)$, so it contains it too. Since
$P^\star\in\DLCN$ was arbitrary, the returned interval contains $\{P(v):P\in\DLCN\}$, which is the
outer-bound claim. No topological hypothesis was used.
\end{proof}

\begin{remark}[The implicit independencies are a relaxation, not an added constraint]
\label{rem:arielindep}
It is sometimes said that ARIEL ``asserts'' pairwise independencies among a factor's boundary
variables that are absent from the LMC --- for the polytree collider factor $\{x_0,x_1,x_2\}$, for
instance, the false $(x_0\indep x_2)$ and $(x_1\indep x_2)$. This would be alarming if those
assertions were \emph{imposed}, since adding independencies shrinks the feasible set and could produce
an unsound inner bound. They are not imposed. The routine that materializes them
(\texttt{\_message\_independencies}, \texttt{ariel.py:652}) is reached only from \texttt{analyze()}
(line~770) and never feeds a solver; the local NLP contains only the simplex, the factor's sentences,
the in-scope LMC equalities and the neighbour boxes, as listed above. The implicit independencies
describe what a marginal-interval message \emph{fails to convey}, i.e.\ an information loss that
\emph{enlarges} the effective feasible set. They therefore threaten tightness, never validity ---
consistent with Lemma~\ref{lem:arielsound}, which shows the true $P^\star$ is always feasible.
\end{remark}

\begin{proposition}[Tightness on a tree of small factors]\label{prop:arieltight}
Suppose (i) the factor graph is a tree; (ii) every factor scope has size $\le2$, so each
factor-to-variable update has at most one other boundary variable; and (iii) the LCN is fully
realizable (Definition~\ref{def:realizable}), and each local NLP is solved to global optimality. Then
the converged intervals equal the exact $\DLCN$ bounds.
\end{proposition}

\begin{proof}
By Lemma~\ref{lem:arielsound} the intervals contain the truth, so it suffices to show each endpoint is
\emph{attained} by some $P\in\DLCN$. Root the factor tree at the queried atom and induct on subtree
depth.

For a leaf factor $f$ with scope $\{n\}$ or $\{m,n\}$ where $m$ has no other factor: the local NLP's
constraints are exactly $f$'s sentences (plus a vacuous box), so an optimizer $\vec p$ is a
distribution on $\mathrm{scope}(f)$ satisfying them, and by (iii) it extends to a member of $\DLCN$ ---
this is where realizability is used: each sentence constrains one family's conditional alone, so the
extension may set the remaining families' conditionals freely without disturbing any sentence, and by
Proposition~\ref{prop:suff} the resulting product lies in $\DLCN$. Hence the endpoint is attained.

For an internal factor $f$ with scope $\{m,n\}$: by (ii) there is exactly one other boundary variable
$m$, and by (i) removing $m$ disconnects $f$'s side of the tree from the rest, so the atoms beyond $m$
form a separate subtree. By induction the incoming interval endpoints on $m$ are attained by
distributions on the far side. Because $m$ is binary, its marginal is the single scalar
$t=P(m{=}1)$, and the incoming interval is $[\lo t,\hi t]$; by induction every value in
$\{\lo t,\hi t\}$ is realized by a far-side member of $\DLCN$, and by continuity (the far-side
constraint set is connected) so is every $t$ in between. The local NLP optimizes $P(n)$ over
distributions on $\{m,n\}$ with $P(m)=t$ for some admissible $t$ and with $f$'s sentences satisfied.
Its optimizer selects a value $t^\star$ and a conditional $P(n\mid m)$; by realizability that
conditional is constrained only by $f$'s own sentences, and by the induction $t^\star$ is realized on
the far side --- independently, since the two sides share only $m$. Composing the two gives a member of
$\DLCN$ attaining the endpoint.
\end{proof}

\begin{remark}[The general claim, and why we state it as cited]\label{rem:arielcited}
ARIEL is asserted to be exact on \emph{all} singly-connected LCNs --- Theorem~1 of Marinescu et al.,
\emph{Approximate Inference in Logical Credal Networks} (IJCAI-2023) --- and the shipped engine does
reproduce the exact bounds on \texttt{chain.lcn}, \texttt{tree.lcn} and \texttt{polytree.lcn} to four
decimals. We state that result with attribution rather than reproving it, because
Proposition~\ref{prop:arieltight}'s hypothesis (ii) is not merely a convenience of the proof.
At a collider factor with scope $\{m_1,m_2,n\}$ --- which occurs on a \emph{polytree}, hence on a
singly-connected graph --- the local NLP constrains $P(m_1)$ and $P(m_2)$ \emph{separately} and then
ranges over \emph{all} joint distributions of $(m_1,m_2)$ with those marginals, including correlated
ones. Two binary variables with fixed marginals admit a one-parameter family of joints, of which the
LMC's co-parent independence $(m_1\indep m_2)$ picks out a single point. So the relaxation is genuine
and does not vanish for binary atoms: single-connectedness alone does not obviously supply tightness,
and we do not claim a proof that it does. What is certain, and proved above, is
Lemma~\ref{lem:arielsound}: ARIEL is always a valid outer bound. This is also the entry the summary
tables need.
\end{remark}

\begin{remark}[Solver and slack caveats]\label{rem:arielsolver}
Two implementation facts qualify Lemma~\ref{lem:arielsound} in practice. The local NLPs are nonconvex
(cleared Type-2 rows and bilinear LMC equalities) and are solved by ipopt, a \emph{local} method: a
local optimum that falls short of the true min/max yields a too-narrow message, which propagates and
can render the result unsound --- the same hazard as \texttt{ExactInference}-local
(Remark~\ref{rem:exactlocal}). Additionally the neighbour box constraints are enforced with slack
variables under a finite penalty, so they can be violated when the penalty is not binding. Both are
reasons to treat ARIEL's output as an approximation rather than a certificate.
\end{remark}

\begin{example}[The loopy boundary, with an uncoupled control]\label{ex:diamond}
Proposition~\ref{prop:arieltight}'s hypothesis (i) is necessary. Take the diamond
$A\to\{B,C\}$, $\{B,C\}\to D$:
\begin{center}\small
$0.4\le P(A)\le0.6$;\quad $P(B\mid A),P(C\mid A)$ transitions;\quad $P(D\mid B,C)$ (four rows);\quad
$\boxed{P(B\wedge C)\le0.25}$.
\end{center}
The boxed cross-family sentence creates the factor $\{B,C\}$ and a second cycle, so the factor graph is
loopy. Measured:
\begin{center}\small
\begin{tabular}{@{}lccc@{}}
\toprule
atom & ARIEL & \texttt{CredalJT} (exact) & verdict\\
\midrule
$A$ & $[0.4000,0.6000]$ & $[0.4000,0.6000]$ & exact\\
$B$ & $[0.3600,\mathbf{0.6400}]$ & $[0.3600,\mathbf{0.5700}]$ & \emph{outer} ($+0.07$)\\
$C$ & $[0.3600,\mathbf{0.6400}]$ & $[0.3600,\mathbf{0.5700}]$ & \emph{outer} ($+0.07$)\\
$D$ & $[0.3160,\mathbf{0.6840}]$ & $[0.3160,\mathbf{0.5940}]$ & \emph{outer} ($+0.09$)\\
\bottomrule
\end{tabular}
\end{center}
The factor $\{B,C\}$ \emph{knows} the cap but can only \emph{report} marginal intervals on $B$ and
$C$; the joint restriction is lost at the message boundary, so the cap never tightens the rest.
\emph{Control:} removing the cap makes ARIEL and \texttt{CredalJT} agree on all four atoms
($A=[0.4,0.6]$, $B=C=[0.36,0.64]$, $D=[0.316,0.684]$), which localizes the looseness to the
cross-family coupling rather than to the loop alone. Consistently with
Lemma~\ref{lem:arielsound}, ARIEL is outer, never inner, throughout.
\end{example}

\begin{remark}[Shipped status: the factor-grouping bug is fixed]\label{rem:arielbug}
An earlier version of \texttt{FactorGraph.\_make\_\allowbreak factor\_nodes} created the first
factor before the grouping loop and then reused the label \texttt{f1} for the next new-scope factor
(the index was incremented after use), silently overwriting --- and dropping the sentences of --- the
first factor. Since sentences are sorted by decreasing scope, the casualty was always the
largest-scope group, so it returned vacuous $[0,1]$ interior marginals on chain/tree/polytree. The fix (each new factor gets a
fresh post-incremented label) is present at \texttt{utils/factor\_graph.py:225--237}. This is worth
recording because the companion note \texttt{ariel\_exactness.tex} documents the post-fix behaviour
without mentioning the bug.
\end{remark}

\section{The exact oracle, and its limits}\label{sec:exact}

\texttt{ExactInference} (\texttt{marginal/exact.py}) solves, per atom, a $\min$ and a $\max$ of $P(a)$
(or the ratio $P(a\mid e)$) over the full $2^n$ joint subject to \emph{all} sentences and \emph{all}
LMC equalities --- i.e.\ directly over $\DLCN$ of \eqref{eq:dlcn}. It is the definitional target, not a
member of the credal-network family.

\begin{proposition}[Certified global; local unreliable]\label{prop:exactinf}
With \texttt{solver="global"} (SCIP spatial branch-and-bound) and a reported gap $\le$
\texttt{gap\_tol} with status \texttt{confirmed}, \texttt{ExactInference} returns the exact $\DLCN$
bounds. With \texttt{solver="local"} (ipopt with an SLSQP fallback) the returned interval is
\emph{unreliable} in the sense of Definition~\ref{def:char}: it may be spuriously inner.
\end{proposition}

\begin{proof}
The optimization model is by construction \eqref{eq:dlcn}, whose feasible set is exactly $\DLCN$;
hence any \emph{globally} optimal value equals the corresponding true bound, and SCIP's spatial
branch-and-bound certifies global optimality up to the reported gap --- when the gap is $0$ and the
status is \texttt{confirmed}, the value is the exact bound. For the local backend, the feasible region
is nonconvex (the LMC equalities \eqref{eq:lmc} and the cleared Type-2 rows \eqref{eq:cleared} are
bilinear), so a local NLP method can converge to a stationary point whose objective is strictly
interior to the true extremum; the multi-restart and SLSQP fallback mitigate but do not certify this,
so no containment in either direction is guaranteed.
\end{proof}

\begin{remark}[\texttt{ExactInference}-local over-tightens]\label{rem:exactlocal}
The failure in Proposition~\ref{prop:exactinf} is observed, not hypothetical: on
\texttt{examples/alarm.lcn} the brute-force and SCIP cross-checks
(\texttt{verify\_bruteforce.py}, \texttt{verify\_scip.py}) put $P(A)=[0.087,0.100]$ and
$P(E)=[0.050,0.074]$, whereas \texttt{solver="local"} reports the spuriously tighter
$P(A)=[0.092,0.100]$, $P(E)=[0.050,0.073]$ --- an inner, unsound interval. Never use the local backend
as a guarantee.
\end{remark}

\begin{remark}[Certified-global does not always terminate]\label{rem:alarmunsolved}
The global backend is not a universal oracle either, and its per-atom \texttt{status} field must be
read. On the five-atom \texttt{alarm.lcn} a run with a $45$\,s per-solve limit returned
$A=[0,0.1]$, $B=[0.1,1.0]$, $E=[0.05,1.0]$, $C=D=[0,1]$ --- but the status shows why: the lower solve
for $A$ and the \emph{upper} solves for $B$ and $E$ came back \texttt{unsolved}, so those endpoints are
fallbacks rather than bounds, while $C$ and $D$ are \texttt{confirmed} genuinely vacuous. These wide
values are consistent with, not a refutation of, the certified $[0.087,0.100]$ etc.\ of
Remark~\ref{rem:exactlocal}; they are simply uncertified. Practical consequences: (i) treat an
\texttt{ExactInference}-global number as exact only when its status is \texttt{confirmed}; (ii) on
instances like this one \texttt{CredalJT} --- which solved every example in this document in seconds
--- is the more trustworthy engine, despite $n$ being small enough that the $2^n$ model is
formulable.
\end{remark}

\section{Summary tables}\label{sec:tables}

\begin{table}[ht]\centering\footnotesize
\caption{Bound character per engine. ``Targets'' is the feasible set the engine optimizes;
``character'' is per Definition~\ref{def:char} \emph{against that set}. The last column gives the
condition for the returned interval to equal the $\DLCN$ truth. Throughout,
\texttt{method="linear"} is assumed (Remark~\ref{rem:d1convex}).}\label{tab:master}
\setlength{\tabcolsep}{4pt}
\begin{tabular}{@{}lllll@{}}
\toprule
Engine & Targets & Character (vs.\ target) & Exact w.r.t.\ $\DLCN$ when & Cost \\
\midrule
\texttt{ExactInference} (global) & $\DLCN$ & exact if \texttt{confirmed} & status \texttt{confirmed} & $2^n$, SCIP \\
\texttt{ExactInference} (local)  & $\DLCN$ & unreliable & --- (may over-tighten) & $2^n$, ipopt \\
\texttt{CredalJT} (D5)           & $\DLCN$ & outer always; exact if H1--H4 & H1--H4 (Thm.~\ref{thm:d5}) & $O(n2^{w+1})$ \\
\midrule
Credal~VE                        & $\DSE$ & exact (uncond., Thm.~\ref{thm:cve}) & Prop.~\ref{prop:suff} holds & per-target elim.\ \\
Credal~CTE ($\epsilon{=}0$)      & $\DSE$ & \textbf{inner} (Prop.~\ref{prop:cte}) & clean s.c.; no cross-family sent. & two-pass tree \\
ApproxLP                         & $\DSE$ & \emph{inner} (Thm.~\ref{thm:alp}) & Prop.~\ref{prop:suff} $+$ no stall & coord.\ descent \\
Interval~BP (\texttt{interval})  & $\DSE$ & outer; exact on tree (Thm.~\ref{thm:ibp}) & tree $+$ Prop.~\ref{prop:suff} & loopy messages \\
Interval~BP (\texttt{variational})& $\DSE$ & \emph{inner} & --- & mean field \\
Credal~IJGP                      & $\DSE$ & outer; exact if $i\ge w$ & $i\ge w$ $+$ Prop.~\ref{prop:suff} & $2^i$/cluster \\
\midrule
ARIEL                            & $\DLCN$ & outer always (Lem.~\ref{lem:arielsound}) & Prop.~\ref{prop:arieltight}; cited Thm.~1 & factor messages \\
SCCP / SCCP-ARIEL                & $\DLCN$ & outer & trivial SCCs & per-SCC NLP \\
\bottomrule
\end{tabular}
\end{table}

\begin{table}[ht]\centering\footnotesize
\caption{Bound character by topology, for unconditional marginals. ``ex'' $=$ exact w.r.t.\ $\DLCN$,
``out'' $=$ valid outer bound, ``in'' $=$ inner (unsound as a guarantee), ``n/f'' $=$ did not finish.
Entries marked $^\dagger$ were measured for this document.}\label{tab:topo}
\setlength{\tabcolsep}{4.5pt}
\begin{tabular}{@{}lcccccc@{}}
\toprule
Topology & Exact-glob & CredalJT & CVE & CTE & ApproxLP & IBP-int \\
\midrule
Markov chain, tree, polytree (realizable)     & ex & ex & ex & ex & ex & ex \\
Chain graph with blocks (\texttt{alarm})      & ex$^\ast$ & ex & ex & ex & ex & ex \\
Cross-family sentence (\texttt{d4\_biting})   & ex$^\dagger$ & ex$^\dagger$ & ex$^\dagger$ & \textbf{in}$^\dagger$ & ex$^\dagger$ & ex$^\dagger$ \\
Non-realizable Type-1 (Ex.~\ref{ex:2node})    & ex$^\dagger$ & ex$^\dagger$ & out$^\dagger$ & out & out$^\dagger$ & out$^\dagger$ \\
Non-realizable Type-2 (Ex.~\ref{ex:collider}) & ex$^\dagger$ & ex$^\dagger$ & out$^\dagger$ & out & out & out \\
Loopy $k$-tree, realizable, $k{=}2$ (Ex.~\ref{ex:loopyexact}) & ex$^\dagger$ & ex$^\dagger$ & \textbf{ex}$^\dagger$ & --- & --- & --- \\
Loopy $k$-tree, realizable, $k{=}3$ (Ex.~\ref{ex:fr5})       & ex$^\dagger$ & ex$^\dagger$ & \textbf{n/f}$^\dagger$ & --- & --- & --- \\
Undirected block (\texttt{smokers})           & ex & \textbf{out} & out & out & in & out \\
\bottomrule
\end{tabular}
\end{table}

\noindent$^\ast$ On \texttt{alarm} the global backend leaves several endpoints \texttt{unsolved}
(Remark~\ref{rem:alarmunsolved}), and the \emph{local} backend over-tightens
(Remark~\ref{rem:exactlocal}); the certified values come from the verifiers.

\begin{remark}[What changed relative to earlier versions of this table]\label{rem:tablechanges}
Three rows were corrected on the basis of measurement. (a) Credal~CTE is no longer listed as ``exact
over $\DSE$'': it is \emph{inner} once a cross-family sentence appears
(Proposition~\ref{prop:cte}). (b) The loopy $k$-tree rows previously read ``out'' for Credal~VE on the
grounds that the class is loopy; at $k{=}2$ Credal~VE is measurably \emph{exact}, and at $k{=}3$ its
problem is non-termination, not looseness (Example~\ref{ex:loopyexact},
Remark~\ref{rem:loopycost}). (c) ApproxLP's character is now stated against $\DSE$, since w.r.t.\
$\DLCN$ it can be outer (Example~\ref{ex:alpouter}).
\end{remark}

\section{Decision guide}\label{sec:guide}

\begin{itemize}
  \item \textbf{Need a guarantee, small $n$:} \texttt{ExactInference} with \texttt{solver="global"}
    --- and check the per-atom \texttt{status}, since \texttt{unsolved} endpoints are not bounds
    (Remark~\ref{rem:alarmunsolved}). Never rely on \texttt{solver="local"}.
  \item \textbf{Need exactness at scale with bounded treewidth:} \texttt{CredalJT} (D5). Exact under
    H1--H4, always a valid outer bound otherwise, and in practice the most trustworthy engine here ---
    including on instances where the $2^n$ oracle fails to certify. Mind
    Remark~\ref{rem:denfloor} for conditional queries.
  \item \textbf{All singleton marginals over $\DSE$, exactly:} Credal~VE, for unconditional queries
    (Theorem~\ref{thm:cve}). Prefer it over Credal~CTE whenever a cross-family sentence is present.
  \item \textbf{A two-sided bracket:} pair ApproxLP (inner over $\DSE$) with Credal~VE (exact over
    $\DSE$) on the same vertices; coinciding endpoints certify the $\DSE$ bound and differing ones
    localize a coordinate-descent stall (Corollary~\ref{cor:bracket}).
  \item \textbf{Fast and singly connected:} Interval~BP \texttt{method="interval"} or ARIEL. Do not use
    Interval~BP for conditional queries with evidence on descendants
    (Remark~\ref{rem:ibpevid}).
  \item \textbf{Suspect a loose bound?} Check realizability first (Definition~\ref{def:realizable}):
    a single non-realizable sentence is the most common cause, and it is a syntactic property you can
    check by inspection via Lemma~\ref{lem:invariance}. Loopiness alone is \emph{not} a reason to
    expect looseness (Example~\ref{ex:loopyexact}).
  \item \textbf{Tightening:} D1--D3 are free build-time knobs for every $\DSE$ engine, but note that D1
    breaks the convexity the vertex machinery assumes (Remark~\ref{rem:d1convex}); D4/D5 address
    Gap~B.
\end{itemize}

\section*{Provenance of the numbers}
\addcontentsline{toc}{section}{Provenance of the numbers}
Every interval reported in Examples~\ref{ex:fr5}, \ref{ex:collider}, \ref{ex:2node},
\ref{ex:loopyexact}, \ref{ex:cveex}, \ref{ex:alpouter}, \ref{ex:ijgp}, \ref{ex:diamond}, in
Proposition~\ref{prop:cte}, and in Remarks~\ref{rem:ibpevid} and~\ref{rem:alarmunsolved} was
produced by running the shipped engines on the stated instance. Values attributed to
\texttt{ExactInference}(\texttt{global}) are reported as exact only where the per-atom status is
\texttt{confirmed} with gap $0$.

Two figures are \emph{not} re-derived here and are cited from the verifiers instead: the certified
\texttt{alarm} marginals of Remark~\ref{rem:exactlocal} (\texttt{verify\_bruteforce.py} and
\texttt{verify\_scip.py}), and the pre-fix Interval~BP numbers of Remark~\ref{rem:ibpbugs}. The
coupled diamond of Example~\ref{ex:diamond} is not shipped as a \texttt{.lcn} file; it was
reconstructed from the description in \texttt{ariel\_exactness.tex} and reproduces that note's
numbers.

\section*{Companion notes}
\addcontentsline{toc}{section}{Companion notes}
Per-engine detail lives in \texttt{cve\_exactness.tex}, \texttt{ccte\_exactness.tex},
\texttt{ibp\_exactness.tex}, \texttt{approxlp\_exactness.tex}, \texttt{ariel\_exactness.tex},
\texttt{cjt\_exactness.tex}, \texttt{strong\_extension\_\allowbreak exactness.tex} and
\texttt{tighter\_approximation.tex} (the D1--D5 taxonomy). Where this document and a companion note
disagree, the discrepancies are deliberate and are flagged at Remarks~\ref{rem:ctecorrection},
\ref{rem:nobiconditional}, \ref{rem:ktreehonest} and~\ref{rem:arielcited}, and in
Lemma~\ref{lem:rip}'s non-uniqueness clause.

\end{document}
